\SourcePage{131}
\SGAchapter{II}{拓扑与层}{Topologies et faisceaux}
\begin{center}J. L. Verdier\end{center}
在定义\Term{topology}{拓扑}与\Term{pretopology}{预拓扑}（第 1 节）和\Term{sheaf-of-sets}{集合层}（第 2 节）之后，第 3 节讨论本讲的\Term{center}{中心}定理（3.4）：\Term{presheaf}{预层}的\Term{associated-sheaf}{伴随层}的存在性。为了涵盖实际中遇到的一切情形，我们对 $\mathcal U$-\Term{site}{景}上的\Term{presheaf}{预层}证明这一\Term{functor}{函子}的存在性（3.0.2）。第 4、5 节从这个定理推出\Term{sheaf}{层}\Term{category}{范畴}中\Term{inductive-limit}{归纳极限}与\Term{projective-limit}{投射极限}的性质。第 6 节定义并研究取值于任意\Term{category}{范畴}的\Term{sheaf}{层}，以便立即着手研究\Term{ring-sheaf}{环层}、\Term{group-sheaf}{群层}、\Term{module-sheaf}{模层}等。

\SGAsection{1}{拓扑、覆盖族、预拓扑}
\begin{SGAresult}{定义}{1.1}
\Term{category}{范畴} $C$ 上的一个\Term{topology}{拓扑}，是对 $C$ 的每个\Term{object}{对象} $X$ 给定一个由 $X$ 的\Term{sieve}{筛}组成的\Term{set}{集合} $J(X)$，并满足下列公理：
\SourcePage{132}
\par T 1)（对\Term{base-change}{基变换}的\Term{stable}{稳定}性。）对 $C$ 的任意\Term{object}{对象} $X$、任意\Term{sieve}{筛} $R\in J(X)$ 以及任意\Term{morphism}{态射} $f:Y\to X$（$Y\in\operatorname{ob}(C)$），$Y$ 的\Term{sieve}{筛} $R\times_XY$ 属于 $J(Y)$。
\par T 2)（\Term{local-scheme}{局部}性。）若 $R$ 与 $R'$ 是 $X$ 的两个\Term{sieve}{筛}，$R\in J(X)$，且对每个 $Y\in\operatorname{ob}(C)$ 和每个\Term{morphism}{态射} $Y\to R$，\Term{sieve}{筛} $R'\times_XY$ 都属于 $J(Y)$，则 $R'$ 属于 $J(X)$。
\par T 3) 对 $C$ 的每个\Term{object}{对象} $X$，$X$ 属于 $J(X)$。
\end{SGAresult}
\textbf{1.1.1。}属于 $J(X)$ 的\Term{sieve}{筛}称为\Term{covering}{覆盖} $X$ 的\Term{sieve}{筛}，也称 $X$ 的\Term{refinement}{加细}。由公理 (T 1)、(T 2)、(T 3) 立即推出，\Term{covering}{覆盖} $X$ 的\Term{sieve}{筛}所成的\Term{set}{集合}对\Term{finite}{有限}交\Term{stable}{稳定}，而包含一个\Term{covering-sieve}{覆盖筛}的\Term{sieve}{筛}也是\Term{covering-sieve}{覆盖筛}。因此，按包含关系排序的\Term{set}{集合} $J(X)$ 是\Term{cofiltered}{余滤过}的（第 I 讲第 8 节）。

\textbf{1.1.2。}设 $C$ 是一个\Term{category}{范畴}，$T$ 与 $T'$ 是 $C$ 上的两个\Term{topology}{拓扑}。若对 $C$ 的每个\Term{object}{对象} $X$，$X$ 对于\Term{topology}{拓扑} $T'$ 的每个\Term{refinement}{加细}也是它对于\Term{topology}{拓扑} $T$ 的\Term{refinement}{加细}，就称\Term{topology}{拓扑} $T$ 比\Term{topology}{拓扑} $T'$ 更细。这就在\Term{topology}{拓扑}的\Term{set}{集合}上定义了一个序结构。

\textbf{1.1.3。}设 $(T_i)_{i\in I}$ 是 $C$ 上的一族\Term{topology}{拓扑}。对每个\Term{object}{对象} $X\in C$，取对于所有\Term{topology}{拓扑} $T_i$ 都\Term{covering}{覆盖} $X$ 的\Term{sieve}{筛}所成的\Term{set}{集合}，这些\Term{set}{集合}满足公理 (T 1)、(T 2)、(T 3)，因而定义一个\Term{topology}{拓扑}：各 $T_i$ 的\Term{intersection-topology}{交拓扑}，即各 $T_i$ 的\Term{infimum}{下确界}。它是在比所有 $T_i$ 都粗的\Term{topology}{拓扑}中最细的。因此，族 $(T_i)_{i\in I}$ 也有\Term{supremum}{上确界}：比每一个 $T_i$ 都细的\Term{topology}{拓扑}的\Term{intersection-topology}{交拓扑}。

\textbf{1.1.4。}特别地，令 $J(X)$ 等于 $X$ 的所有\Term{sieve}{筛}组成的\Term{set}{集合}，就得到一个比 $C$ 上任何\Term{topology}{拓扑}都细的\Term{topology}{拓扑}，我们称之为\Term{discrete-topology}{离散拓扑}。
存在一个比 $C$ 上任何\Term{topology}{拓扑}都粗的\Term{topology}{拓扑}，即对每个 $X\in C$ 都有 $J(X)=\{X\}$ 的\Term{topology}{拓扑}。它称为\Term{coarse-topology}{粗拓扑}或\Term{chaotic-topology}{混沌拓扑}。

\textbf{1.1.5。}配备一个\Term{topology}{拓扑}的\Term{category}{范畴} $C$ 称为\Term{site}{景}。\Term{category}{范畴} $C$ 称为这个\Term{site}{景}的底\Term{category}{范畴}。

\begin{SGAresult}{定义}{1.2}
设 $C$ 是一个\Term{site}{景}，$X$ 是 $C$ 的一个\Term{object}{对象}。\Term{morphism}{态射}族 $(f_\alpha:X_\alpha\to X)$，$\alpha\in A$，称为\Term{covering-family}{覆盖族}，如果由族 $f_\alpha$ 生成的\Term{sieve}{筛}（第 I 讲 4.3.3）是\Term{covering}{覆盖} $X$ 的\Term{sieve}{筛}。
\end{SGAresult}
\textbf{1.1.6。}设 $C$ 是一个\Term{category}{范畴}。对 $C$ 的每个\Term{object}{对象} $X$，给定一些以 $X$ 为靶的\Term{morphism}{态射}族所成的\Term{set}{集合}。则使这些给定的族成为\Term{covering-family}{覆盖族}的\Term{topology}{拓扑}中存在一个最粗的\Term{topology}{拓扑} $T$，即所有这些\Term{topology}{拓扑}的交（1.1.3）。称这个\Term{topology}{拓扑}为由所给各个\Term{morphism}{态射}族的\Term{set}{集合}生成的\Term{topology}{拓扑}。
\SourcePage{133}
一般说来，难以描述这个\Term{topology}{拓扑}的所有\Term{covering-sieve}{覆盖筛}。不过，在下述情形中情况比较好：
\begin{SGAresult}{定义}{1.3}
设 $C$ 是一个\Term{category}{范畴}。$C$ 上的一个\Term{pretopology}{预拓扑}，是对每个\Term{object}{对象} $X\in C$ 给定一个由以 $X$ 为靶的\Term{morphism}{态射}族组成的\Term{set}{集合} $\operatorname{Cov}(X)$，并满足下列公理：
\par PT 0) 对任意\Term{object}{对象} $X\in C$，$\operatorname{Cov}(X)$ 中各\Term{morphism}{态射}族的\Term{morphism}{态射}都是\Term{quarrable-morphism}{可作基变换}的。（回忆，称\Term{morphism}{态射} $Y\to X$ \Term{quarrable-morphism}{可作基变换}，是指对每个\Term{morphism}{态射} $Z\to X$，\Term{fiber-product}{纤维积} $Z\times_XY$ \Term{representable}{可表示}。）
\par PT 1) 对任意\Term{object}{对象} $X\in C$、属于 $\operatorname{Cov}(X)$ 的任意族 $(X_\alpha\to X)_{\alpha\in A}$，以及任意\Term{morphism}{态射} $Y\to X$（$Y\in\operatorname{ob}(C)$），族 $(X_\alpha\times_XY\to Y)_{\alpha\in A}$ 属于 $\operatorname{Cov}(Y)$。（对\Term{base-change}{基变换}\Term{stable}{稳定}。）
\par PT 2) 若 $(X_\alpha\to X)_{\alpha\in A}$ 属于 $\operatorname{Cov}(X)$，且对每个 $\alpha\in A$，$(X_{\beta_\alpha}\to X_\alpha)_{\beta_\alpha\in B_\alpha}$ 属于 $\operatorname{Cov}(X_\alpha)$，则族 $(X_\gamma\to X)_{\gamma\in\coprod_{\alpha\in A}B_\alpha}$ 属于 $\operatorname{Cov}(X)$；这里，对 $\gamma=(\alpha,\beta_\alpha)$，$\alpha\in A$、$\beta_\alpha\in B_\alpha$，\Term{morphism}{态射} $X_\gamma\to X$ 是复合\Term{morphism}{态射} $X_\gamma=X_{\beta_\alpha}\to X_\alpha\to X$。（对复合\Term{stable}{稳定}。）
\par PT 3) 族 $X\xrightarrow{\mathrm{id}_X}X$ 属于 $\operatorname{Cov}(X)$。
\end{SGAresult}
\textbf{1.3.1。}由 1.1.6 的定义，对 $C$ 上给定的\Term{pretopology}{预拓扑}，可以考虑由它生成的 $C$ 上的\Term{topology}{拓扑}。注意，若 $C$ 是一个\Term{fiber-product}{纤维积}都\Term{representable}{可表示}的\Term{category}{范畴}，则 $C$ 的任何\Term{topology}{拓扑} $T$ 都可由一个\Term{pretopology}{预拓扑}定义：令 $\operatorname{Cov}(X)$ 由对于\Term{topology}{拓扑} $T$ \Term{covering}{覆盖} $X$ 的一切族组成即可。

\begin{SGAresult}{命题}{1.4}
设 $C$ 是一个\Term{category}{范畴}，$E$ 是 $C$ 上的\Term{pretopology}{预拓扑}，$T$ 是\Term{pretopology}{预拓扑} $E$ 定义的\Term{topology}{拓扑}（1.1.6），$X$ 是 $C$ 的\Term{object}{对象}。以 $J_E(X)$ 表示由\Term{pretopology}{预拓扑}中的族生成的 $X$ 的\Term{sieve}{筛}所成的\Term{set}{集合}，以 $J_T(X)$ 表示 $X$ 对于\Term{topology}{拓扑} $T$ 的\Term{refinement}{加细}所成的\Term{set}{集合}。则 $J_E(X)$ 在 $J_T(X)$ 中\Term{cofinal}{共尾}。换言之，$X$ 的\Term{sieve}{筛} $R$ 属于 $J_T(X)$ 的充要条件是存在\Term{sieve}{筛} $R'\in J_E(X)$ 使 $R'\subset R$。
\end{SGAresult}
\begin{proof}
对任意\Term{object}{对象} $X\in C$，以 $J'(X)$ 表示包含某个 $J_E(X)$ 的\Term{sieve}{筛}的那些 $X$ 的\Term{sieve}{筛}所成的\Term{set}{集合}。显然有 $J'(X)\subset J_T(X)$。要证 $J'(X)=J_T(X)$，只须证明各 $J'(X)$ 定义 $C$ 上的一个\Term{topology}{拓扑}。$J'(X)$ 显然满足公理 (T 1)、(T 3)，只剩下 (T 2)。为此，只须证明：若 $R'$ 是 $R\in J_E(X)$ 的子\Term{sieve}{筛}，且对任意 $Y\to R$，$Y$ 的\Term{sieve}{筛} $R'\times_XY$ 属于 $J'(X)$，则 $R'$ 属于 $J'(X)$。而\Term{sieve}{筛} $R$ 由一个属于 $\operatorname{Cov}(X)$ 的族 $(X_\alpha\to X)$ 生成；对每个 $\alpha$，\Term{sieve}{筛} $R'\times_XX_\alpha$ 包含一个由属于 $\operatorname{Cov}(X_\alpha)$ 的族 $(X_{\beta_\alpha}\to X_\alpha)$ 生成的\Term{sieve}{筛}。因此，$R'$ 包含由族 $(X_{\beta_\alpha}\to X)$ 生成的\Term{sieve}{筛}。由公理 (PT 2)，$R'$ 包含一个 $J_E(X)$ 中的\Term{sieve}{筛}，故属于 $J'(X)$。
\end{proof}

\SourcePage{134}
\SGAsection{2}{集合层}
\begin{SGAresult}{定义}{2.1}
设 $C$ 是底\Term{category}{范畴}为 $\mathcal U$-\Term{category}{范畴}的\Term{site}{景}。取值于 $\mathcal U\text{-}\mathrm{Ens}$ 的\Term{presheaf}{预层} $F$ 称为\Term{separated}{分离}的（相应地，称为\Term{sheaf}{层}），如果对\Term{covering}{覆盖} $C$ 的\Term{object}{对象} $X$ 的每个\Term{sieve}{筛} $R$，映射
\[\operatorname{Hom}_{\widehat C}(X,F)\longrightarrow\operatorname{Hom}_{\widehat C}(R,F)\]
是\Term{injectivity}{单射}（相应地，是\Term{bijection}{双射}）。$\widehat C$ 中以\Term{sheaf}{层}为\Term{object}{对象}的\Term{full-subcategory}{满子范畴}，称为 $C$ 上的\Term{sheaf-of-sets}{集合层}\Term{category}{范畴}，通常记为 $\widetilde C$\footnote{或者记为 $C^{\sim}$，视打字机的心情而定。}。在不致混淆时，我们简称 $C$ 上的\Term{sheaf}{层}\Term{category}{范畴}；有时则更明确地称为取值于 $\mathcal U\text{-}\mathrm{Ens}$ 的\Term{sheaf}{层}\Term{category}{范畴}，或 $\mathcal U$-\Term{sheaf}{层}\Term{category}{范畴}。
\end{SGAresult}
\begin{SGAresult}{命题}{2.2}
设 $C$ 是一个 $\mathcal U$-\Term{category}{范畴}，$\mathscr F=(F_i)_{i\in I}$ 是 $C$ 上的一族 $\mathcal U$-\Term{presheaf}{预层}。对每个\Term{object}{对象} $X\in C$，令 $J_{\mathscr F}(X)$ 为满足下述条件的\Term{sieve}{筛} $R\to X$ 的\Term{set}{集合}：对 $C$ 中以 $X$ 为靶的每个\Term{morphism}{态射} $Y\to X$，\Term{sieve}{筛} $R\times_XY$ 使得映射
\[\operatorname{Hom}_{\widehat C}(Y,F_i)\longrightarrow\operatorname{Hom}_{\widehat C}(R\times_XY,F_i)\]
对所有 $i\in I$ 都是\Term{bijection}{双射}（相应地，是\Term{injectivity}{单射}）。则这些\Term{set}{集合} $J_{\mathscr F}(X)$ 定义 $C$ 上的一个\Term{topology}{拓扑}，且它是使每个 $F_i$ 成为\Term{sheaf}{层}（相应地，\Term{separated-presheaf}{分离预层}）的最细\Term{topology}{拓扑}。
\end{SGAresult}
\begin{proof}
$J_{\mathscr F}(X)$ 显然满足公理 (T 1)、(T 3)，只须证明 (T 2)。为此，只须证明：
\par 1) 若 $R'\hookrightarrow R\hookrightarrow X$ 是 $X$ 的两个\Term{sieve}{筛}，$R\in J_{\mathscr F}(X)$，且对每个 $Y\to R$（$Y$ 是 $C$ 的\Term{object}{对象}），\Term{sieve}{筛} $R'\times_XY$ 属于 $J_{\mathscr F}(Y)$，则 $R'$ 属于 $J_{\mathscr F}(X)$。
\par 2) 若 $R'\hookrightarrow R\hookrightarrow X$ 是 $X$ 的两个\Term{sieve}{筛}，且 $R'\in J_{\mathscr F}(X)$，则\Term{sieve}{筛} $R$ 属于 $J_{\mathscr F}(X)$。
注意，在情形 2) 中，对任意 $Y\to R$（$Y$ 是 $C$ 的\Term{object}{对象}），\Term{sieve}{筛} $R'\times_XY$ 属于 $J_{\mathscr F}(Y)$。在情形 1)、2) 中，$R$ 是 $R$ 上方的 $C$ 的\Term{object}{对象} $Y$ 的\Term{inductive-limit}{归纳极限}（第 I 讲 3.4）。$\widehat C$ 中的\Term{inductive-limit}{归纳极限}是\Term{universal}{普适}的（第 I 讲 3.3），所以在两种情形中，$R'$ 都是各 $R'\times_XY$ 的\Term{inductive-limit}{归纳极限}。而这两种情形下 $R'\times_XY$ 都属于 $J_{\mathscr F}(Y)$。对 $R$ 上方的 $C$ 的\Term{object}{对象} $Y$ 取\Term{inductive-limit}{归纳极限}，得到映射
\[\operatorname{Hom}_{\widehat C}(R,F_i)\longrightarrow\operatorname{Hom}_{\widehat C}(R',F_i)\]
对每个 $i\in I$ 都是\Term{bijection}{双射}（相应地，\Term{injectivity}{单射}）。因而映射
\[\operatorname{Hom}_{\widehat C}(X,F_i)\longrightarrow\operatorname{Hom}_{\widehat C}(R',F_i),\qquad
\operatorname{Hom}_{\widehat C}(X,F_i)\longrightarrow\operatorname{Hom}_{\widehat C}(R,F_i)\]
\SourcePage{135}
在情形 1)、2) 中都是\Term{bijection}{双射}（相应地，\Term{injectivity}{单射}）。此外，假设 1)、2) 显然对任意\Term{base-change}{基变换} $Y\to X$（$Y$ 是 $C$ 的\Term{object}{对象}）\Term{stable}{稳定}。因此，对 $X$ 上方的 $C$ 的每个\Term{object}{对象} $Y$ 及每个 $i\in I$，映射
\[\begin{split}
\operatorname{Hom}_{\widehat C}(Y,F_i)&\longrightarrow\operatorname{Hom}_{\widehat C}(R\times_XY,F_i),\\
\operatorname{Hom}_{\widehat C}(Y,F_i)&\longrightarrow\operatorname{Hom}_{\widehat C}(R'\times_XY,F_i)
\end{split}\]
在两种情形中都是\Term{bijection}{双射}（相应地，\Term{injectivity}{单射}）；这证明 $R'$、$R$ 属于 $J_{\mathscr F}(X)$。
\end{proof}
\begin{SGAresult}{推论}{2.3}
设 $C$ 是一个 $\mathcal U$-\Term{category}{范畴}，对每个 $X\in C$ 给定一个由 $X$ 的\Term{sieve}{筛}组成的\Term{set}{集合} $K(X)$。假定各 $K(X)$ 满足 (1.1) 的公理 (T 1)。\Term{presheaf}{预层} $F$ 对于由各 $K(X)$ 生成的\Term{topology}{拓扑}（1.1.6）是\Term{sheaf}{层}（相应地，是\Term{separated-presheaf}{分离预层}）的充要条件是，对 $C$ 的每个\Term{object}{对象} $X$ 和每个\Term{sieve}{筛} $R\in K(X)$，映射
\[\operatorname{Hom}_{\widehat C}(X,F)\longrightarrow\operatorname{Hom}_{\widehat C}(R,F)\]
是\Term{bijection}{双射}（相应地，是\Term{injectivity}{单射}）。
\end{SGAresult}
\begin{SGAresult}{推论}{2.4}
特别地，设 $C$ 是配备一个\Term{pretopology}{预拓扑}的 $\mathcal U$-\Term{category}{范畴}。\Term{presheaf}{预层} $F$ 是\Term{sheaf}{层}（相应地，是\Term{separated-presheaf}{分离预层}）的充要条件是，对 $C$ 的每个\Term{object}{对象} $X$ 及属于 $\operatorname{Cov}(X)$ 的每个族 $(X_\alpha\to X)$，\Term{set}{集合}图
\[F(X)\longrightarrow\prod_{\alpha\in A}F(X_\alpha)\rightrightarrows
\prod_{(\alpha,\beta)\in A\times A}F(X_\alpha\times_XX_\beta)\]
是\Term{exact}{正合}的（相应地，映射 $F(X)\longrightarrow\prod_{\alpha\in A}F(X_\alpha)$ 是\Term{injectivity}{单射}）。
\end{SGAresult}
\begin{proof}应用 2.3，然后应用第 I 讲 3.5 和第 I 讲 2.12。\end{proof}
推论 2.4 给出了 [1] 中的定义。
\begin{SGAresult}{定义}{2.5}
设 $C$ 是一个 $\mathcal U$-\Term{category}{范畴}。称使\Term{representable-functor}{可表示函子}成为\Term{sheaf}{层}的最细\Term{topology}{拓扑}（2.2）为 $C$ 的\Term{canonical-topology}{典范拓扑}。对于\Term{canonical-topology}{典范拓扑}\Term{covering}{覆盖} $X$ 的\Term{sieve}{筛}，称为\Term{universally-strictly-epimorphic-sieve}{普适严格满筛}。对于\Term{canonical-topology}{典范拓扑}的\Term{covering-family}{覆盖族}，称为\Term{universally-strictly-epimorphic-family}{普适严格满态射族}。若\Term{covering-family}{覆盖族}中的\Term{morphism}{态射}还都是\Term{quarrable-morphism}{可作基变换}的，则称这个族为\Term{universally-effective-epimorphic-family}{普适有效满态射族} [2]。
\end{SGAresult}
\textbf{注记 2.5.1。}在迄今所用的几乎所有\Term{site}{景}中，\Term{topology}{拓扑}都比\Term{canonical-topology}{典范拓扑}粗；换言之，$C$ 上的\Term{representable-functor}{可表示函子}都是\Term{sheaf}{层}，即 $C$ 的\Term{covering-family}{覆盖族}都是
\SourcePage{136}
\Term{universally-strictly-epimorphic-family}{普适严格满态射族}。这条规则唯一的例外是第 5 节研究的\Term{site}{景} $\widehat C$（它的\Term{topology}{拓扑}比\Term{canonical-topology}{典范拓扑}细，而且通常严格更细）。
\begin{SGAresult}{命题}{2.6}
\Term{sieve}{筛} $R\hookrightarrow X$ 是\Term{universally-strictly-epimorphic-sieve}{普适严格满筛}的充要条件是，对 $C$ 中 $X$ 上方的每个\Term{object}{对象} $Y\to X$ 和 $C$ 的每个\Term{object}{对象} $Z$，映射
\[\operatorname{Hom}_C(Y,Z)\longrightarrow
\varprojlim_{C/(Y\times_XR)}\operatorname{Hom}_C(\,\cdot\,,Z)\]
是\Term{bijection}{双射}。
\end{SGAresult}
\begin{proof}应用 2.2 和第 I 讲 5.3 即得。\end{proof}
\textbf{注记 2.7。}1) 命题 2.6 给出了\Term{category}{范畴} $C$ 的\Term{universally-strictly-epimorphic-sieve}{普适严格满筛}的刻画，与\Term{presheaf}{预层}取值所在的\Term{universe}{宇宙}无关，唯一条件是 $C$ 的\Term{morphism}{态射}\Term{set}{集合} $\operatorname{Hom}(X,Y)$ 属于这个\Term{universe}{宇宙}。因此，无须指定\Term{universe}{宇宙}，即可对任意\Term{category}{范畴} $C$ 定义\Term{canonical-topology}{典范拓扑}。

2) 设 $C$ 是底\Term{category}{范畴}为 $\mathcal U$-\Term{category}{范畴}的\Term{site}{景}，$F$ 是一个 $\mathcal U$-\Term{presheaf}{预层}，$\mathcal V$ 是包含 $\mathcal U$ 的\Term{universe}{宇宙}。$C$ 的底\Term{category}{范畴}是一个 $\mathcal V$-\Term{category}{范畴}，且 $F$ 可以看成一个 $\mathcal V$-\Term{presheaf}{预层}。$\mathcal U$-\Term{presheaf}{预层} $F$ 是\Term{sheaf}{层}（相应地，是\Term{separated-presheaf}{分离预层}），当且仅当 $\mathcal V$-\Term{presheaf}{预层} $F$ 是\Term{sheaf}{层}（相应地，是\Term{separated-presheaf}{分离预层}）。换言之，\Term{presheaf}{预层} $F$ 成为\Term{sheaf}{层}（相应地，\Term{separated-presheaf}{分离预层}）的条件，在刚才明确的意义下，与\Term{presheaf}{预层} $F$ 取值所在的\Term{universe}{宇宙}无关。特别地，设 $C$ 是一个\Term{site}{景}，$F$ 是一个 $\mathcal U$-\Term{presheaf}{预层}；若存在包含 $\mathcal U$ 的\Term{universe}{宇宙} $\mathcal V$，使\Term{category}{范畴} $C$ 是 $\mathcal V$-\Term{category}{范畴}且 $F$ 是 $\mathcal V$-\Term{sheaf}{层}（相应地，\Term{separated}{分离}的 $\mathcal V$-\Term{presheaf}{预层}），就说 $F$ 是\Term{sheaf}{层}（相应地，\Term{separated-presheaf}{分离预层}）。这一性质与\Term{universe}{宇宙} $\mathcal V$ 无关。

\SGAsection{3}{预层的伴随层}
\begin{SGAresult}{定义}{3.0.1}
设 $C$ 是一个\Term{site}{景}。称 $C$ 的\Term{object}{对象}组成的\Term{set}{集合} $G$ 为 $C$ 的\Term{topologically-generating-family}{拓扑生成族}（不致混淆时也简称\Term{generating-family}{生成族}），如果 $C$ 的每个\Term{object}{对象}都是一个源属于 $G$ 的 $C$ 中\Term{morphism}{态射}\Term{covering-family}{覆盖族}（1.2）的靶。
\end{SGAresult}
\begin{SGAresult}{定义}{3.0.2}
设 $\mathcal U$ 是一个\Term{universe}{宇宙}。称底\Term{category}{范畴}是 $\mathcal U$-\Term{category}{范畴}（第 I 讲 1.1）且有一个小\Term{topologically-generating-family}{拓扑生成族}的\Term{site}{景} $C$ 为 $\mathcal U$-\Term{site}{景}。设 $C$ 是一个 $\mathcal U$-\Term{category}{范畴}；称使 $C$ 成为 $\mathcal U$-\Term{site}{景}的 $C$ 上\Term{topology}{拓扑}为 $\mathcal U$-\Term{topology}{拓扑}。若\Term{site}{景} $C$ 的底\Term{category}{范畴}是小的（第 I 讲 1.0），就说 $C$ 是 $\mathcal U$-小的，也常不甚严格地简称为小\Term{site}{景}。
\end{SGAresult}
\SourcePage{137}
\textbf{3.0.3。}由定义立即得出，比一个 $\mathcal U$-\Term{topology}{拓扑}更细的任何\Term{topology}{拓扑}仍是 $\mathcal U$-\Term{topology}{拓扑}，而小\Term{site}{景}总是 $\mathcal U$-\Term{site}{景}。
\begin{SGAresult}{命题}{3.0.4}
设 $C$ 是一个 $\mathcal U$-\Term{site}{景}，$G$ 是 $C$ 的一个小\Term{topologically-generating-family}{拓扑生成族}。对每个 $X\in\operatorname{ob}C$，以 $J_G(X)$ 表示由形如 $(Y_\alpha\xrightarrow{u_\alpha}X)$、$\alpha\in A$、$Y_\alpha\in G$ 的\Term{morphism}{态射}族生成的\Term{covering}{覆盖} $X$ 的\Term{sieve}{筛}所成的\Term{set}{集合}。
\par 1) \Term{set}{集合} $J_G(X)$ 是小的。
\par 2) \Term{set}{集合} $J_G(X)$ 在按包含关系排序的所有\Term{covering}{覆盖} $X$ 的\Term{sieve}{筛}的\Term{set}{集合} $J(X)$ 中\Term{cofinal}{共尾}。
\par 3) 对任意 $R\in J_G(X)$，存在一个小\Term{epimorphic-family}{满态射族}（第 I 讲 10.2）$(u_\alpha:Y\to R)$，$\alpha\in A$，其中 $Y\in G$。
\end{SGAresult}
\begin{proof}
1) 令 $A(X)=\coprod_{Y\in G}\operatorname{Hom}(Y,X)$。\Term{set}{集合} $A(X)$ 是小的（第 I 讲第 0 节），且 $\operatorname{card}(J_G(X))\leq 2^{\operatorname{card}(A(X))}$，所以 $J_G(X)$ 是小的。

2) 设 $R\in J(X)$。令 $A(R)=\coprod_{Y\in G}\operatorname{Hom}(Y,R)$，令 $R'$ 为族 $(Y\xrightarrow{u}R\hookrightarrow X)$、$u\in A(R)$ 生成的 $X$ 的\Term{sieve}{筛}。有 $R'\subset R$，只须证明 $R'$ \Term{covering}{覆盖} $X$。由 1.1 的公理 (T 2)，只须证明对每个\Term{morphism}{态射} $Z\to R$，$Z\in\operatorname{ob}C$，$Z$ 的\Term{sieve}{筛} $R'\times_XZ$ 是\Term{covering-sieve}{覆盖筛}。而 $R'\times_XZ$ 包含由所有 $Y\to Z$、$Y\in G$ 组成的族生成的\Term{sieve}{筛}；按假设，这个族是\Term{covering-family}{覆盖族}。因此，$Z$ 的\Term{sieve}{筛} $R'\times_XZ$ 是\Term{covering-sieve}{覆盖筛}（1.1 的公理 T 2）。

3) 设 $R\in J_G(X)$。按假设，族 $(Y\xrightarrow{u}R)$、$u\in A(R)=\coprod_{Y\in G}\operatorname{Hom}(Y,R)$ 是\Term{epimorphic-family}{满态射族}。对每个 $Y\in\operatorname{ob}C$，$\operatorname{Hom}(Y,R)$ 包含于\Term{small-set}{小集合} $\operatorname{Hom}(Y,X)$，所以 $A(R)$ 是小的。
\end{proof}
\textbf{3.0.5。}设 $C$ 是一个 $\mathcal U$-\Term{site}{景}，$\mathcal V$ 是满足 $C\in\mathcal V$、$\mathcal U\subset\mathcal V$ 的\Term{universe}{宇宙}，$G$ 是 $C$ 的一个 $\mathcal U$-小\Term{topologically-generating-family}{拓扑生成族}。$C$ 上的 $\mathcal U$-\Term{presheaf-of-sets}{集合预层}\Term{category}{范畴} $\widehat C$ 是 $\mathcal V$-\Term{category}{范畴}（第 I 讲 1.1.1）。设 $X$ 是 $C$ 的\Term{object}{对象}。\Term{covering}{覆盖} $X$ 的\Term{sieve}{筛}的\Term{set}{集合} $J(X)$ 是 $\mathcal V$-小的，且按包含关系排序是\Term{cofiltered}{余滤过}的（1.1.1）。故对每个 $\mathcal U$-\Term{presheaf}{预层} $F$，\Term{inductive-limit}{归纳极限} $\varinjlim_{J(X)}\operatorname{Hom}_{\widehat C}(R,F)$ 可由 $\mathcal V$ 的一个元素表示（第 I 讲 2.4.1）。又由 3.0.4 的 3)，对每个 $R\in J_G(X)$，$\operatorname{Hom}_{\widehat C}(R,F)$ 是 $\mathcal U$-小的；而 $J_G(X)$ 是 $J(X)$ 中的一个 $\mathcal U$-小\Term{cofinal}{共尾}\Term{set}{集合}（同上），所以由第 I 讲 2.4.2，$\varinjlim_{J(X)}\operatorname{Hom}_{\widehat C}(R,F)$ 是 $\mathcal U$-小的。对每个 $F$ 和每个 $X$，选择一个表示此\Term{inductive-limit}{归纳极限}的 $\mathcal U$ 的元素，令
\[LF(X)=\varinjlim_{J(X)}\operatorname{Hom}_{\widehat C}(\,\cdot\,,F).\]
设 $g:Y\to X$ 是 $C$ 的\Term{morphism}{态射}。\Term{base-change}{基变换}\Term{functor}{函子} $g^*:J(X)\to J(Y)$ 定义一个映射
\[LF(g):LF(X)\longrightarrow LF(Y),\]
\SourcePage{138}
使 $X\mapsto LF(X)$ 成为 $C$ 上的\Term{presheaf}{预层}。
由于\Term{morphism}{态射} $\mathrm{id}_X:X\to X$ 是 $J(X)$ 的一个元素，对 $C$ 的每个\Term{object}{对象} $X$，有一个映射
\[\ell(F)(X):F(X)\longrightarrow LF(X),\]
它显然定义一个\Term{functor}{函子}\Term{morphism}{态射} $\ell(F):F\to LF$。又显然，$F\mapsto LF$ 对 $F$ 是\Term{functor}{函子}，而各\Term{morphism}{态射} $\ell(F)$ 定义一个\Term{morphism}{态射}
\[\ell:\mathrm{Id}\longrightarrow L.\]
最后，设 $R\hookrightarrow X$ 是 $X$ 的\Term{refinement}{加细}。$LF(X)$ 的定义和第 I 讲 1.4 给出一个映射
\[Z_R:\operatorname{Hom}_{\widehat C}(R,F)\longrightarrow\operatorname{Hom}_{\widehat C}(X,LF),\]
且对 $C$ 中任意\Term{morphism}{态射} $Y\xrightarrow{g}X$，由\Term{functor}{函子} $LF$ 的定义，下面的图交换：
\[
\begin{tikzcd}[column sep=large]
\operatorname{Hom}_{\widehat C}(R,F)\arrow[r,"Z_R"]\arrow[d]&\operatorname{Hom}_{\widehat C}(X,LF)\arrow[d]\\
\operatorname{Hom}_{\widehat C}(R\times_XY,F)\arrow[r,"Z_{R\times_XY}"]&\operatorname{Hom}_{\widehat C}(Y,LF)
\end{tikzcd}\tag{*}
\]
（竖\Term{arrow}{箭头}为显然的映射。）现在照录 [2] 的论证。
\begin{SGAresult}{引理}{3.1}
1) 对每个\Term{refinement}{加细} $i_R:R\hookrightarrow X$ 和每个 $u:R\to F$，图
\[
\begin{tikzcd}
F\arrow[r,"\ell(F)"]&LF\\
R\arrow[u,"u"]\arrow[r,hook,"i_R"']&X\arrow[u,"Z_R(u)"']
\end{tikzcd}\tag{**}
\]
交换。
\par 2) 对每个\Term{morphism}{态射} $v:X\to LF$，存在 $X$ 的一个\Term{refinement}{加细} $R$ 以及\Term{morphism}{态射} $u:R\to F$，使 $Z_R(u)=v$。
\par 3) 设 $Y$ 是 $C$ 的\Term{object}{对象}，$u,v:Y\rightrightarrows F$ 是满足 $\ell(F)\circ u=\ell(F)\circ v$ 的两个\Term{morphism}{态射}。则\Term{morphism}{态射}对 $(u,v)$ 的\Term{equalizer}{等化子}是 $Y$ 的\Term{refinement}{加细}。
\SourcePage{139}
\par 4) 设 $R$、$R'$ 是 $X$ 的两个\Term{refinement}{加细}，$u:R\to F$、$u':R'\to F$ 是两个\Term{morphism}{态射}。$Z_R(u)=Z_{R'}(u')$ 的充要条件是，$u$ 和 $u'$ 在某个\Term{refinement}{加细} $R''\hookrightarrow R\times_XR'$ 上相同。
\end{SGAresult}
\begin{proof}
唯一非\Term{trivial-bundle}{平凡}的断言是 1)。须证 $Z_R(u)\circ i_R=\ell(F)\circ u$。为此，只须证明这两个\Term{morphism}{态射}与每个\Term{morphism}{态射} $Y\xrightarrow{g}R$（$Y$ 是 $C$ 的\Term{object}{对象}）的复合相同（第 I 讲 3.4）。令 $f=i_R\circ g$，考虑\Term{fiber-product}{纤维积} $R\times_XY$：
\[
\begin{tikzcd}[column sep=large,row sep=large]
F\arrow[r,"\ell(F)"]&LF\\
R\arrow[u,"u"]\arrow[r,hook,"i_R"']&X\arrow[u,"Z_R(u)"']\\
R\times_XY\arrow[u,"g'"]\arrow[r,"i'"',"\sim"]&Y\arrow[u,"f"']\arrow[ul,"g"']
\end{tikzcd}
\]
上图中的 $i'$ 是具有\Term{section}{截面}的\Term{monomorphism}{单态射}，故是\Term{isomorphism}{同构}（第 I 讲 1.3 的 3)）。由\Term{morphism}{态射} $\ell(F)$ 的定义，\Term{morphism}{态射} $Z_{R\times_XY}(u\circ g')$ 等于 $\ell(F)\circ u\circ g$。另一方面，图 (*) 的交换性给出 $Z_{R\times_XY}(u\circ g')=Z_R(u)\circ f$，于是确有 $\ell(F)\circ u\circ g=Z_R(u)\circ i_R\circ g$。
\end{proof}
\begin{SGAresult}{命题}{3.2}
1) \Term{functor}{函子} $L$ \Term{left-exact}{左正合}（第 I 讲 2.3.2）。
\par 2) 对每个\Term{presheaf}{预层} $F$，$LF$ 是\Term{separated-presheaf}{分离预层}。
\par 3) \Term{presheaf}{预层} $F$ \Term{separated}{分离}，当且仅当\Term{morphism}{态射} $\ell(F):F\to LF$ 是\Term{monomorphism}{单态射}。此时\Term{presheaf}{预层} $LF$ 是\Term{sheaf}{层}。
\par 4) 下列性质\Term{equivalence}{等价}：
\par (i) $\ell(F):F\to LF$ 是\Term{isomorphism}{同构}；
\par (ii) $F$ 是\Term{sheaf}{层}。
\end{SGAresult}
\begin{proof}
1) 只须证明，对 $C$ 的每个\Term{object}{对象} $X$，\Term{functor}{函子} $F\mapsto LF(X)$ 与\Term{finite-projective-limit}{有限投射极限}交换（第 I 讲 3.1）。由\Term{projective-limit}{投射极限}的定义，\Term{functor}{函子}
\[F\longmapsto\operatorname{Hom}_{\widehat C}(R,F),\qquad R\hookrightarrow X\in\operatorname{ob}(J(X))\]
与\Term{projective-limit}{投射极限}交换，而\Term{inductive-limit}{归纳极限} $\varinjlim_{J(X)}$ 与\Term{finite-projective-limit}{有限投射极限}交换，因为 $J(X)$ 是\Term{cofiltered}{余滤过}\Term{set}{集合}（第 I 讲 2.7）。

2) 设 $X$ 是 $C$ 的\Term{object}{对象}，$f,g:X\rightrightarrows LF$ 是在 $X$ 的某个\Term{refinement}{加细} $R\hookrightarrow X$ 上相同的两个\Term{morphism}{态射}。由 3.1 的 2)，存在\Term{covering-sieve}{覆盖筛} $R'\hookrightarrow X$，
\SourcePage{140}
总可设它包含于 $R$，以及两个\Term{morphism}{态射} $u,v:R'\rightrightarrows F$，使 $Z_{R'}(u)=f$、$Z_{R'}(v)=g$。由 3.1 的 1)，有 $\ell(F)\circ u=\ell(F)\circ v$。故由 3.1 的 4)，$u$ 与 $v$ 在某个\Term{refinement}{加细} $R''\hookrightarrow R'$ 上相同。以 $w$ 表示 $u$ 在 $R''$ 上的\Term{restriction}{限制}，则 $Z_{R''}(w)=Z_{R'}(u)=Z_{R'}(v)$，故 $f=g$。所以\Term{presheaf}{预层} $LF$ \Term{separated}{分离}。

3) 若 $F$ \Term{separated}{分离}，则\Term{morphism}{态射} $\ell(F)$ 是\Term{monomorphism}{单态射}，因为\Term{monomorphism}{单态射}的\Term{filtered-inductive-limit}{滤过归纳极限}是\Term{monomorphism}{单态射}。若 $\ell(F)$ 是\Term{monomorphism}{单态射}，则\Term{presheaf}{预层} $F$ 是一个\Term{separated-presheaf}{分离预层}的\Term{subpresheaf}{子预层}，因而\Term{separated}{分离}。证明此时 $LF$ 是\Term{sheaf}{层}。设 $i:R\hookrightarrow X$ 是\Term{covering}{覆盖} $C$ 的一个\Term{object}{对象} $X$ 的\Term{sieve}{筛}，$u:R\to LF$ 是\Term{morphism}{态射}。只须证明 $u$ 经过 $X$ \Term{resolution}{分解}。令 $R'=F\times_{LF}R$，$v=Z_{R'}(\mathrm{pr}_1)$：
\[
\begin{tikzcd}[column sep=large,row sep=large]
F\arrow[r,"\ell(F)"]&LF\\
R'\arrow[u,"\mathrm{pr}_1"]\arrow[r,"\mathrm{pr}_2"]&R\arrow[u,"u"]\arrow[r,"i"']&X\arrow[ul,"v"']\\
R''\arrow[u,"p_1"]\arrow[r,"p_2"']&Y\arrow[u,"m"']
\end{tikzcd}
\]
为证 $u=v\circ i$，只须证明对每个\Term{morphism}{态射} $m:Y\to R$（$Y$ 是 $C$ 的\Term{object}{对象}），$v\circ i\circ m=u\circ m$（第 I 讲 3.4）。令 $R''=Y\times_RR'$，以 $p_1$、$p_2$ 表示\Term{projection}{投影}。\Term{projection}{投影} $R''\overset{p_2}{\hookrightarrow}Y$ 是\Term{monomorphism}{单态射}，并使 $R''$ 成为\Term{covering}{覆盖} $Y$ 的\Term{sieve}{筛}（3.1 的 2)）。有 $v\circ i\circ\mathrm{pr}_2=\ell(F)\circ\mathrm{pr}_1$（3.1 的 1)），从而 $v\circ i\circ\mathrm{pr}_2=u\circ m\circ p_2$。由此推出 $v\circ i\circ\mathrm{pr}_2\circ p_1=u\circ\mathrm{pr}_2\circ p_1$，即 $v\circ i\circ m\circ p_2=u\circ m\circ p_2$。由于\Term{presheaf}{预层} $LF$ \Term{separated}{分离}，得到 $v\circ i\circ m=u\circ m$。

4) 显然。
\end{proof}
\textbf{注记 3.3。}设 $J'(X)$ 是 $J(X)$ 的一个\Term{cofinal}{共尾}\Term{subset}{子集}，则
\[LF(X)=\varinjlim_{J'(X)}\operatorname{Hom}_{\widehat C}(\,\cdot\,,F).\]
特别地，若 $C$ 的\Term{topology}{拓扑}由\Term{pretopology}{预拓扑} $X\mapsto\operatorname{Cov}(X)$ 定义（1.1.3），则\Term{functor}{函子} $L$ 可以用 $\operatorname{Cov}(X)$ 中的\Term{covering-family}{覆盖族}描述（第 I 讲 2.12 和第 I 讲 3.5）。把公式展开，就得到 [ 中的构造。
\begin{SGAresult}{定理}{3.4}
设 $C$ 是一个 $\mathcal U$-\Term{site}{景}。从\Term{sheaf}{层}到\Term{presheaf}{预层}的\Term{inclusion-functor}{包含函子} $i:\widetilde C\hookrightarrow\widehat C$ 有一个\Term{left-exact}{左正合}的\Term{left-adjoint}{左伴随} $a$（第 I 讲 2.3.2）：
\[\begin{tikzcd}\widetilde C\arrow[r,hook,shift right=0.6ex,"i"']&\widehat C\arrow[l,shift right=0.6ex,"a"']\end{tikzcd}\]
\SourcePage{141}
\Term{functor}{函子} $i\circ a$ \Term{canonical}{典范}\Term{isomorphism}{同构}于\Term{functor}{函子} $L\circ L$（参见 3.0.5）。对每个\Term{presheaf}{预层} $F$，\Term{adjunction-morphism}{伴随态射} $F\to i\circ(F)$ 是由上述\Term{isomorphism}{同构}从\Term{morphism}{态射} $\ell(LF)\ell(F):F\to L\circ L(F)$ 得到的。
\end{SGAresult}
\begin{SGAresult}{定义}{3.5}
\Term{sheaf}{层} $F$ 称为\Term{presheaf}{预层} $F$ 的\Term{associated-sheaf}{伴随层}。
\end{SGAresult}
定理 3.4 由命题 3.2 立即得出。
\begin{SGAresult}{命题}{3.6}
设 $C$ 是一个 $\mathcal U$-\Term{site}{景}，$\mathcal V\supset\mathcal U$ 是一个\Term{universe}{宇宙}。以 $\widehat C_{\mathcal U}$、$\widetilde C_{\mathcal U}$（相应地，$\widehat C_{\mathcal V}$、$\widetilde C_{\mathcal V}$）表示 $\mathcal U$-\Term{presheaf}{预层}和 $\mathcal U$-\Term{sheaf}{层}（相应地，$\mathcal V$-\Term{presheaf}{预层}和 $\mathcal V$-\Term{sheaf}{层}）的\Term{category}{范畴}，以 $a_{\mathcal U}:\widehat C_{\mathcal U}\to\widetilde C_{\mathcal U}$（相应地，$a_{\mathcal V}:\widehat C_{\mathcal V}\to\widetilde C_{\mathcal V}$）表示相应的“\Term{associated-sheaf}{伴随层}”\Term{functor}{函子}。图
\[\begin{tikzcd}
\widehat C_{\mathcal U}\arrow[r,"a_{\mathcal U}"]\arrow[d,hook]&\widetilde C_{\mathcal U}\arrow[d,hook]\\
\widehat C_{\mathcal V}\arrow[r,"a_{\mathcal V}"']&\widetilde C_{\mathcal V}
\end{tikzcd}\]
在\Term{canonical}{典范}\Term{isomorphism}{同构}意义下交换，其中竖直\Term{functor}{函子}为\Term{canonical}{典范}包含。
\end{SGAresult}
\begin{proof}由\Term{functor}{函子} $a_{\mathcal U}$、$a_{\mathcal V}$ 的构造（3.4 和 3.0.5）得出。\end{proof}

\SGAsection{4}{层范畴的正合性质}
\Term{sheaf}{层}\Term{category}{范畴}的\Term{exact}{正合}性质通过定理 3.4 由\Term{presheaf}{预层}\Term{category}{范畴}的\Term{exact}{正合}性质\Term{derivation}{导出}。本节选取其中最有用的一些典型命题，具体说明这一思路。
\begin{SGAresult}{定理}{4.1}
设 $C$ 是一个 $\mathcal U$-\Term{site}{景}，$\widetilde C$ 是\Term{sheaf}{层}\Term{category}{范畴}，$a:\widehat C\to\widetilde C$ 是\Term{associated-sheaf-functor}{伴随层函子}，$i:\widetilde C\to\widehat C$ 是\Term{inclusion-functor}{包含函子}。
\par 1) \Term{functor}{函子} $a$ 与\Term{inductive-limit}{归纳极限}交换，且是\Term{exact}{正合}的。
\par 2) $\widetilde C$ 中的 $\mathcal U$-\Term{inductive-limit}{归纳极限}\Term{representable}{可表示}。对任意小\Term{category}{范畴} $I$ 和任意\Term{functor}{函子} $E:I\to\widetilde C$，\Term{canonical}{典范}\Term{morphism}{态射}
\[\varinjlim_I E\longrightarrow a(\varinjlim_I i\circ E)\]
是\Term{isomorphism}{同构}。
\par 3) $\widetilde C$ 中的 $\mathcal U$-\Term{projective-limit}{投射极限}\Term{representable}{可表示}。对 $C$ 的每个\Term{object}{对象} $X$，$\widetilde C$ 上的\Term{functor}{函子} $F\mapsto F(X)$ 与\Term{projective-limit}{投射极限}交换，即\Term{inclusion-functor}{包含函子} $i:\widetilde C\to\widehat C$ 与\Term{projective-limit}{投射极限}交换。
\end{SGAresult}
\begin{proof}这些性质基本上由定理 3.4 和第 I 讲 2.11 得出。\end{proof}
\SourcePage{142}
因此，在\Term{sheaf}{层}\Term{category}{范畴}中，以 $\mathcal U$ 的元素为指标的积、\Term{fiber-product}{纤维积}、以 $\mathcal U$ 的元素为指标的和、\Term{pushout}{纤维和}、核、\Term{cokernel}{余核}、像和\Term{coimage}{余像}都\Term{representable}{可表示}。
\begin{SGAresult}{推论}{4.1.1}
设 $C$ 是一个 $\mathcal U$-\Term{site}{景}，$F$ 是 $C$ 上的\Term{sheaf-of-sets}{集合层}。\Term{canonical-homomorphism}{典范同态}
\[\varinjlim_{C/F}X\longrightarrow F\]
是\Term{isomorphism}{同构}。
\end{SGAresult}
\begin{proof}由第 I 讲 3.4 及 $a$ 与\Term{inductive-limit}{归纳极限}交换得出。\end{proof}
\begin{SGAresult}{命题}{4.2}
$\widetilde C$ 中同时为\Term{epimorphism}{满态射}与\Term{monomorphism}{单态射}的每个\Term{morphism}{态射}都是\Term{isomorphism}{同构}。
\end{SGAresult}
\begin{proof}
设 $u:G\to H$ 是 $\widetilde C$ 中同时为\Term{epimorphism}{满态射}和\Term{monomorphism}{单态射}的\Term{morphism}{态射}。首先注意，$u$ 是\Term{presheaf}{预层}的\Term{monomorphism}{单态射}（4.1 的 1)）。在\Term{presheaf}{预层}\Term{category}{范畴}中构造\Term{pushout}{纤维和} $H\coprod_GH=K$ 及两个\Term{canonical}{典范}\Term{morphism}{态射} $i_1,i_2:H\rightrightarrows K$。由于 $u$ 是\Term{presheaf}{预层}的\Term{monomorphism}{单态射}，立即可知下图既是\Term{cartesian}{笛卡尔}的又是\Term{cocartesian}{余笛卡尔}的（第 I 讲 10.1）：
\[\begin{tikzcd}G\arrow[r,"u"]\arrow[d,"u"']&H\arrow[d,"i_2"]\\H\arrow[r,"i_1"']&K\end{tikzcd}\]
施以“\Term{associated-sheaf}{伴随层}”\Term{functor}{函子}，得到\Term{sheaf}{层}\Term{category}{范畴}中既是\Term{cartesian}{笛卡尔}又是\Term{cocartesian}{余笛卡尔}的图（4.1 的 1)）：
\[\begin{tikzcd}G\arrow[r,"u"]\arrow[d,"u"']&H\arrow[d,"ai_2"]\\H\arrow[r,"ai_1"']&aK\end{tikzcd}\tag{*}\]
由于 $u$ 是\Term{sheaf}{层}的\Term{epimorphism}{满态射}，\Term{morphism}{态射} $i_1$ 是\Term{isomorphism}{同构}；又因图 (*) 是\Term{cartesian}{笛卡尔}的，\Term{morphism}{态射} $u$ 是\Term{isomorphism}{同构}。
\end{proof}
\begin{SGAresult}{命题}{4.3}
1) $\widetilde C$ 中\Term{representable}{可表示}的\Term{inductive-limit}{归纳极限}是\Term{universal}{普适}的（第 I 讲 2.5）。
\par 2) 每个\Term{epimorphic-family}{满态射族}（第 I 讲 10.2）都是\Term{universally-effective-epimorphic-family}{普适有效满态射族}（2.6）。
\par 3) 每个\Term{equivalence-relation}{等价关系}都是\Term{universally-effective}{普适有效}的（第 I 讲 10.6）。
\SourcePage{143}
\par 4) $\mathcal U$-\Term{filtered-inductive-limit}{滤过归纳极限}与\Term{finite-projective-limit}{有限投射极限}交换（第 I 讲 2.6）。
\end{SGAresult}
\begin{proof}
断言 1) 至 4) 在\Term{set}{集合}\Term{category}{范畴}中成立，因而在\Term{presheaf}{预层}\Term{category}{范畴}中成立（第 I 讲 3.1）。于是由\Term{presheaf}{预层}的相应断言和 4.1 立即得到 1)、4)。证明 3)。设 $p_1,p_2:R\rightrightarrows X$ 是 $\widetilde C$ 的一个\Term{equivalence-relation}{等价关系}。图 $R\overset{p_1}{\underset{p_2}{\rightrightarrows}}X$ 也是\Term{presheaf}{预层}的\Term{equivalence-relation}{等价关系}（4.1.1)）。所以存在\Term{presheaf}{预层}\Term{morphism}{态射} $u:X\to Y$，使图
\[\begin{tikzcd}R\arrow[r,"p_2"]\arrow[d,"p_1"']&X\arrow[d,"u"]\\X\arrow[r,"u"']&Y\end{tikzcd}\]
在\Term{presheaf}{预层}\Term{category}{范畴}中既是\Term{cartesian}{笛卡尔}又是\Term{cocartesian}{余笛卡尔}的。施以“\Term{associated-sheaf}{伴随层}”\Term{functor}{函子}，得到\Term{sheaf}{层}\Term{category}{范畴}中既是\Term{cartesian}{笛卡尔}又是\Term{cocartesian}{余笛卡尔}的图（4.1 的 1)）。所以 $R\overset{p_1}{\underset{p_2}{\rightrightarrows}}X$ 是\Term{effective-equivalence-relation}{有效等价关系}。由上所述，每个\Term{equivalence-relation}{等价关系}都\Term{effective}{有效}，故每个\Term{equivalence-relation}{等价关系}都\Term{universally-effective}{普适有效}。
证明 2)。设 $(u_i:X_i\to X)$，$i\in I$，是 $\widetilde C$ 的一个\Term{epimorphic-family}{满态射族}。由第 II 讲 2.6 及第 I 讲 2.12，只须证明下列断言：
\par a) 对每个\Term{sheaf}{层}\Term{morphism}{态射} $v:Y\to X$，\Term{morphism}{态射}族 $\mathrm{pr}_{2,i}:X_i\times_XY\to Y$，$i\in I$，是 $\widetilde C$ 的\Term{epimorphic-family}{满态射族}。
\par b) 对每个\Term{sheaf}{层} $Z$，\Term{set}{集合}图
\[\operatorname{Hom}(X,Z)\longrightarrow\prod_{i\in I}\operatorname{Hom}(X_i,Z)\rightrightarrows
\prod_{(i,k)\in I\times I}\operatorname{Hom}(X_i\times_XX_k,Z)\]
\Term{exact}{正合}。
令 $X'$ 为各\Term{morphism}{态射} $u_i$ 在\Term{presheaf}{预层}意义下的像的并，$j:X'\hookrightarrow X$ 为\Term{canonical}{典范}\Term{embedding}{嵌入}。对每个 $i\in I$，\Term{morphism}{态射} $u_i:X_i\to X$ \Term{resolution}{分解}为\Term{morphism}{态射} $u'_i:X_i\to X'$ 与 $j$ 的复合；族 $u'_i:X_i\to X'$，$i\in I$，是 $\widehat C$ 的\Term{epimorphic-family}{满态射族}。由于 $j$ 是\Term{monomorphism}{单态射}，\Term{morphism}{态射} $aj:aX'\to X$ 是\Term{sheaf}{层}的\Term{monomorphism}{单态射}（4.1）。由于对每个 $i\in I$ 有 $u_i=aj\,au'_i$，$aj$ 是\Term{sheaf}{层}的\Term{epimorphism}{满态射}，故 $aj$ 是\Term{isomorphism}{同构}（4.2）。设 $v:Y\to X$ 是\Term{sheaf}{层}\Term{morphism}{态射}。\Term{base-change}{基变换}给出\Term{morphism}{态射} $j_v:X'\times_XY\to Y$ 及各\Term{morphism}{态射} $u'_{i,v}:X_i\times_XY\to X'\times_XY$。由于 $a$ 与\Term{fiber-product}{纤维积}交换，$aj_v$ 是\Term{isomorphism}{同构}。由于 $\widehat C$ 的\Term{epimorphic-family}{满态射族}在\Term{base-change}{基变换}后仍保持这一性质，族 $u'_{i,v}$，$i\in I$，在 $\widehat C$ 中是\Term{epimorphic-family}{满态射族}。由于 $a$ 与\Term{inductive-limit}{归纳极限}交换，族
\SourcePage{144}
$au'_{i,v}:X_i\times_XY\to a(X'\times_XY)$，$i\in I$，在 $\widetilde C$ 中是\Term{epimorphic-family}{满态射族}，得到 a)。证明 b)。设 $Z$ 是\Term{sheaf}{层}。由于 $aj$ 是\Term{isomorphism}{同构}，映射
\[\operatorname{Hom}(X,Z)\longrightarrow\operatorname{Hom}(X',Z)\]
是\Term{bijection}{双射}。由于各 $u'_i$，$i\in I$，组成 $\widehat C$ 的\Term{epimorphic-family}{满态射族}，\Term{set}{集合}图
\[\operatorname{Hom}(X',Z)\longrightarrow\prod_{i\in I}\operatorname{Hom}(X_i,Z)\rightrightarrows
\prod_{(i,k)\in I\times I}\operatorname{Hom}(X_i\times_{X'}X_k,Z)\]
\Term{exact}{正合}。最后，$X'$ 是 $X$ 的\Term{subpresheaf}{子预层}，故对 $(i,k)\in I\times I$，\Term{presheaf}{预层} $X_i\times_{X'}X_k$ \Term{canonical}{典范}\Term{isomorphism}{同构}于 $X_i\times_XX_k$，得到 b)。
\end{proof}
\textbf{注记 4.3.1。}1) 命题 4.3 的 2) 从形式上给出 4.2 的第二个证明：显然，在任何\Term{category}{范畴}中，同时为\Term{universally-effective-epimorphism}{普适有效满态射}和\Term{monomorphism}{单态射}的\Term{morphism}{态射}必为\Term{isomorphism}{同构}。

2) 由上述命题可知，\Term{sheaf}{层}\Term{category}{范畴}中的每个\Term{morphism}{态射}都唯一\Term{resolution}{分解}为一个\Term{effective-epimorphism}{有效满态射}与一个\Term{effective-monomorphism}{有效单态射}的复合。这个性质是在非\Term{additive-category}{可加范畴}情形下对\Term{abelian-category}{阿贝尔范畴}公理 (AB2)（$\operatorname{coim}\simeq\operatorname{im}$）的自然推广 [Tohoku]。

\textbf{4.4.0。}设 $C$ 是一个 $\mathcal U$-\Term{site}{景}。将\Term{functor}{函子} $h:C\to\widehat C$（第 I 讲 1.3.1）与\Term{associated-sheaf-functor}{伴随层函子}复合，得到\Term{functor}{函子}
\[\varepsilon_C:C\longrightarrow\widetilde C,\]
称为从 $C$ 到 $\widetilde C$ 的\Term{canonical-functor}{典范函子}，后面将经常使用。\Term{functor}{函子} $\varepsilon_C$ 与\Term{finite-projective-limit}{有限投射极限}交换。当 $C$ 的\Term{topology}{拓扑}比\Term{canonical-topology}{典范拓扑}粗时，$\varepsilon_C$ \Term{fully-faithful}{全忠实}，且与\Term{projective-limit}{投射极限}交换；而且此时即使 $C$ 不一定有一个小的——而且，即使 $C$ 不一定有一个小\Term{topologically-generating-family}{拓扑生成族}，它也有定义。我们不详细研究 $\varepsilon_C$ 对\Term{inductive-limit}{归纳极限}的性质（参见 [SGA 3 IV]），不过会用到以下命题。
\begin{SGAresult}{定理}{4.4}
设 $C$ 是一个 $\mathcal U$-\Term{site}{景}，$(u_i:X_i\to X)$，$i\in I$，是 $C$ 中以 $X$ 为靶的一族\Term{morphism}{态射}。下列条件\Term{equivalence}{等价}\footnote{当\Term{site}{景} $C$ 不一定有一个小\Term{topologically-generating-family}{拓扑生成族}，而 $C$ 的\Term{topology}{拓扑}比\Term{canonical-topology}{典范拓扑}粗时，仍有 ii)$\Rightarrow$i)（参见 4.4 的证明）。}：
\par i) 族 $(\varepsilon_Cu_i:\varepsilon_CX_i\to\varepsilon_CX)$，$i\in I$，是 $\widetilde C$ 的\Term{epimorphic-family}{满态射族}（第 I 讲 10.2）。
\par ii) 族 $(u_i:X_i\to X)$，$i\in I$，是 $C$ 的\Term{covering-family}{覆盖族}（1.2）。
\end{SGAresult}
\SourcePage{145}
\begin{proof}
ii)$\Rightarrow$i)。设 $R\hookrightarrow X$ 是由各 $u_i$，$i\in I$，生成的\Term{sieve}{筛}（第 I 讲 4.3.3），$u'_i:X_i\to R$ 是由 $u_i$ 诱导的\Term{morphism}{态射}。族 $u'_i$，$i\in I$，是 $\widehat C$ 的\Term{epimorphic-family}{满态射族}，而 $R$ 是\Term{covering}{覆盖} $X$ 的\Term{sieve}{筛}。所以，对每个\Term{sheaf}{层} $F$，映射
\[\operatorname{Hom}(X,F)\longrightarrow\operatorname{Hom}(R,F)\]
是\Term{bijection}{双射}，而映射
\[\operatorname{Hom}(R,F)\longrightarrow\prod_i\operatorname{Hom}(X_i,F)\]
是\Term{injectivity}{单射}。因此映射
\[\operatorname{Hom}(X,F)\longrightarrow\prod_i\operatorname{Hom}(X_i,F)\]
是\Term{injectivity}{单射}，从而映射
\[\operatorname{Hom}(X,F)\longrightarrow\prod_i\operatorname{Hom}(X_i,F)\]
是\Term{injectivity}{单射}，得到 i)。

i)$\Rightarrow$ii)。沿用前面的记号，以 $J_R:R\hookrightarrow X$ 表示各 $u_i$，$i\in I$，生成的\Term{sieve}{筛}到 $X$ 的\Term{canonical}{典范}\Term{embedding}{嵌入}。由 i)，\Term{sheaf}{层}\Term{morphism}{态射} $aJ_R:aR\to aX$ 同时是\Term{sheaf}{层}的\Term{epimorphism}{满态射}和\Term{monomorphism}{单态射}，故为\Term{sheaf}{层}的\Term{isomorphism}{同构}（4.2）。用第 3 节的记号，有\Term{commutative-diagram}{交换图}
\[\begin{tikzcd}[column sep=large]
X\arrow[r,"\ell(X)"]&LX\arrow[r,"\ell(LX)"]&aX\\
R\arrow[u,hook,"J_R"]\arrow[r,"\ell(R)"']&LR\arrow[u,hook,"LJ_R"]\arrow[r,"\ell(LR)"']&aR\arrow[u,"aJ_R"',"\sim"]
\end{tikzcd}\]
首先，由 3.1 的 2)，存在\Term{covering-sieve}{覆盖筛} $J_{S_1}:S_1\to X$ 和\Term{morphism}{态射} $u_1:S_1\to LR$，使得
\[\ell(LX)\circ\ell(X)\circ J_{S_1}=aJ_R\circ\ell(LR)\circ u_1
=\ell(LX)\circ LJ_R\circ u_1.\tag{4.6.1}\]
令 $m=\ell(X)\circ J_{S_1}$、$n=LJ_R\circ u_1$。$m$、$n$ 是从 $S_1$ 到 $LX$ 的\Term{morphism}{态射}。令 $S_2\hookrightarrow S_1$ 为\Term{arrow}{箭头}对 $(m,n)$ 的\Term{equalizer}{等化子}。对 $C$ 的每个\Term{object}{对象} $Y$ 及每个\Term{morphism}{态射} $\alpha:Y\to S_1$，由 4.6.1 有 $\ell(LX)\circ m\circ\alpha=\ell(LX)\circ n\circ\alpha$。由 3.1 的 3)，\Term{arrow}{箭头}对 $(m\circ\alpha,n\circ\alpha)$ 的\Term{equalizer}{等化子} $S_2\times_{S_1}Y$ 是\Term{covering}{覆盖} $Y$ 的\Term{sieve}{筛}。故由\Term{topology}{拓扑}公理 (T 2)，$J_{S_2}:S_2\hookrightarrow X$ 是\Term{covering}{覆盖} $X$ 的\Term{sieve}{筛}。于是有一个\Term{morphism}{态射} $u_2:S_2\to LR$ 和\Term{commutative-diagram}{交换图}：
\SourcePage{146}
\[\begin{tikzcd}[column sep=large,row sep=large]
S_2\arrow[rr,hook,"J_{S_2}"]\arrow[dd,"u_2"']&&X\arrow[dd,"\ell(X)"]\\
&R\arrow[ur,hook,"J_R"]\arrow[dl,"\ell(R)"]&\\
LR\arrow[rr,"LJ_R"']&&LX
\end{tikzcd}\tag{4.6.2}\]
设 $Y$ 是 $C$ 的\Term{object}{对象}，$\beta:Y\to S_2$ 是\Term{morphism}{态射}。由 3.1 的 1)，存在\Term{covering-sieve}{覆盖筛} $J_{Q_1}:Q_1\hookrightarrow Y$ 及\Term{morphism}{态射} $v_1:Q_1\to R$，使 $u_2\circ\beta\circ J_{Q_1}=\ell(R)\circ v_1$。令 $f=J_{S_2}\circ\beta\circ J_{Q_1}$、$g=J_R\circ v_1$。令 $Q_2$ 为\Term{morphism}{态射}对 $(f,g):Q_1\rightrightarrows X$ 的\Term{equalizer}{等化子}，$J_{Q_2}:Q_2\hookrightarrow Y$ 为\Term{canonical}{典范}\Term{embedding}{嵌入}。对 $C$ 的每个\Term{object}{对象} $Z$ 及每个\Term{morphism}{态射} $\gamma:Z\to Q_1$，由 (4.6.2) 的交换性有
\[\begin{split}
\ell(X)\circ f\circ\gamma
&=\ell(X)\circ J_{S_2}\circ\beta\circ J_{Q_1}\circ\gamma
=LJ_R\circ u_2\circ\beta\circ J_{Q_1}\circ\gamma\\
&=LJ_R\circ\ell(R)\circ v_1\circ\gamma
=\ell(X)\circ J_R\circ v_1\circ\gamma
=\ell(X)\circ g\circ\gamma.
\end{split}\]
由 3.1 的 2)，\Term{morphism}{态射}对 $(f\circ\gamma,g\circ\gamma)$ 的\Term{equalizer}{等化子} $Q_2\times_{Q_1}Z$ 是\Term{covering}{覆盖} $Z$ 的\Term{sieve}{筛}。由\Term{topology}{拓扑}公理 (T 2)，$J_{Q_2}:Q_2\hookrightarrow Y$ 是\Term{covering}{覆盖} $Y$ 的\Term{sieve}{筛}。因此，对每个\Term{morphism}{态射} $\beta:Y\to S$，存在\Term{covering-sieve}{覆盖筛} $J_{Q_2}:Q_2\hookrightarrow Y$ 和\Term{morphism}{态射} $v_2:Q_2\to R$，使下图交换：
\[\begin{tikzcd}
Y\arrow[r,"\beta"]&S_2\arrow[r,"J_{S_2}"]&X\\
Q_2\arrow[u,hook,"J_{Q_2}"]\arrow[r,"V_2"']&R\arrow[ur,hook,"J_R"']&
\end{tikzcd}\]
以 $S_2\cap R$ 表示 $S_2$ 与 $R$ 在 $X$ 上的\Term{fiber-product}{纤维积}。$Y$ 的\Term{sieve}{筛} $(S_2\cap R)\times_{S_2}Y$ 包含\Term{sieve}{筛} $Q_2$，因此是\Term{covering}{覆盖} $Y$ 的\Term{sieve}{筛}。再由\Term{topology}{拓扑}公理 (T 2)，$S_2\cap R$ 是\Term{covering}{覆盖} $X$ 的\Term{sieve}{筛}，因而包含它的\Term{sieve}{筛} $R$ 也是\Term{covering}{覆盖} $X$ 的\Term{sieve}{筛}。
\end{proof}
\begin{SGAresult}{推论}{4.4.4}
设 $T$ 是 $\mathcal U$-\Term{site}{景} $e$ 的\Term{topology}{拓扑}，$T'$（相应地，$T''$）是 $e$ 上使 $\widetilde C$ 的\Term{object}{对象}成为\Term{sheaf}{层}（相应地，\Term{presheaf}{预层}）的最细\Term{topology}{拓扑}（2.2）。则 $T=T'=T''$。
\end{SGAresult}
事实上，显然有 $T\leq T'\leq T''$，而 4.4 和 2.2 蕴涵 $T''\leq T$。

\SourcePage{147}
\begin{SGAresult}{定义}{4.5}
\Term{category}{范畴} $A$ 的\Term{initial-object}{始对象}，是表示空\Term{inductive-limit}{归纳极限}的\Term{object}{对象} $\phi_A$，即对每个 $X\in\operatorname{ob}A$，恰有一个\Term{arrow}{箭头} $\phi_A\to X$。若\Term{object}{对象} $\phi_A$ 是\Term{initial-object}{始对象}，且每个以 $\phi_A$ 为靶的\Term{morphism}{态射}都是\Term{isomorphism}{同构}，就称它为\Term{strict-initial-object}{严格始对象}。设 $(S_i)$，$i\in I$，是\Term{category}{范畴} $A$ 的一族\Term{object}{对象}，假设和 $s=\coprod_{i\in I}S_i$ \Term{representable}{可表示}。称和 $S$ 为\Term{disjoint-coproduct}{不交和}，如果结构\Term{morphism}{态射} $S_i\to S$ \Term{quarrable-morphism}{可作基变换}且为\Term{monomorphism}{单态射}，并且对每对 $i,j$，$i\ne j$，积 $S_i\times_SS_j$ 是 $A$ 的\Term{initial-object}{始对象}。若和 $S$ 是\Term{disjoint-coproduct}{不交和}，且在任意\Term{base-change}{基变换} $T\to S$ 后仍为\Term{disjoint-coproduct}{不交和}，就称它为\Term{universal-disjoint-coproduct}{普适不交和}；于是对每对 $i,j$，$i\ne j$，\Term{object}{对象} $S_i\times_SS_j$ 都是 $A$ 的\Term{strict-initial-object}{严格始对象}。
\end{SGAresult}
\textbf{例 4.5.1。}在\Term{set}{集合}\Term{category}{范畴}中，\Term{direct-sum}{直和}是不交且\Term{universal}{普适}的。故在任何\Term{presheaf-of-sets}{集合预层}\Term{category}{范畴}（第 I 讲 3.1）中也如此，因而在任何 $\mathcal U$-\Term{site}{景}上的\Term{sheaf-of-sets}{集合层}\Term{category}{范畴}中也如此（4.1）；特别地，$\widetilde C$ 的\Term{initial-object}{始对象}是严格的。
\begin{SGAresult}{命题}{4.6}
设 $C$ 是一个 $\mathcal U$-\Term{category}{范畴}，$(s_i:X_i\to X)$，$i\in I$，是 $C$ 中的一族\Term{morphism}{态射}。对 $C$ 上每个 $\mathcal U$-\Term{topology}{拓扑} $T$（3.0.2），以 $\widetilde C_T$ 表示相应的\Term{sheaf}{层}\Term{category}{范畴}，以 $\varepsilon_T:C\to\widetilde C_T$ 表示相应\Term{functor}{函子}。
\par 1) 设 $T$ 是一个 $\mathcal U$-\Term{topology}{拓扑}，且
\[\coprod_{i\in I}\varepsilon_T(X_i)\xrightarrow{(\varepsilon_T(s_i))}\varepsilon_T(X)\]
是\Term{isomorphism}{同构}。则对每个比 $T$ 细的\Term{topology}{拓扑} $T'$，\Term{morphism}{态射}
\[\coprod_{i\in I}\varepsilon_{T'}(X_i)\xrightarrow{(\varepsilon_{T'}(s_i))}\varepsilon_{T'}(X)\]
是\Term{isomorphism}{同构}。
\par 2) 设 $T$ 是 $C$ 上的一个 $\mathcal U$-\Term{topology}{拓扑}。下列性质 (i)、(ii) \Term{equivalence}{等价}：
\par i) a) 族 $(s_i:X_i\to X)$，$i\in I$，对于 $T$ 是\Term{covering-family}{覆盖族}。
\par b) 对每个 $i\in I$，\Term{presheaf}{预层}的\Term{diagonal-morphism}{对角态射} $\Delta_i:X_i\hookrightarrow X_i\times_XX_i$ 经“\Term{associated-sheaf}{伴随层}”\Term{functor}{函子}（对于 $T$）变为\Term{isomorphism}{同构}（当各 $s_i$ 都是\Term{monomorphism}{单态射}时就有这一性质）。
\par c) 对 $I$ 的每对元素 $(i,j)$，$i\ne j$，\Term{presheaf}{预层} $X_i\times_XX_j$ 经“\Term{associated-sheaf}{伴随层}”\Term{functor}{函子}（对于 $T$）变为 $\widetilde C_T$ 的\Term{initial-object}{始对象}。
\par ii) $\varepsilon_T(X)$ 是各 $\varepsilon_T(X_i)$ 的和。
\end{SGAresult}
\SourcePage{148}
\begin{proof}
1) 对于 $T'$ 的\Term{sheaf}{层} $F$ 也是对于 $T$ 的\Term{sheaf}{层}。所以对每个对于 $T'$ 的\Term{sheaf}{层} $F$，\Term{morphism}{态射}
\[\operatorname{Hom}_{\widehat C}(X,F)\longrightarrow\prod_{i\in I}\operatorname{Hom}_{\widehat C}(X_i,F)\]
是\Term{isomorphism}{同构}，由此得到断言。

2) 按定义，性质 ii) 成立当且仅当\Term{presheaf}{预层}\Term{morphism}{态射}
\[\Phi=(h(s_i),i\in I):\coprod_i h(X_i)\longrightarrow h(X)\]
经“\Term{associated-sheaf}{伴随层}”\Term{functor}{函子}变为\Term{isomorphism}{同构}，即当且仅当（4.2）$a(\Phi)$ 同时是\Term{sheaf}{层}的\Term{epimorphism}{满态射}和\Term{monomorphism}{单态射}。由 4.4，\Term{morphism}{态射} $a(\Phi)$ 是\Term{epimorphism}{满态射}当且仅当有性质 a)。\Term{diagonal-morphism}{对角态射}
\[\Delta:\coprod_i h(X_i)\longrightarrow
\left(\coprod_i h(X_i)\right)\times_{h(X)}\left(\coprod_i h(X_i)\right)\]
是一族\Term{morphism}{态射} $\Delta_{i,j}$，$(i,j)\in i\times I$，的\Term{direct-sum}{直和}，其中：
\par $\beta$) 当 $i=j$ 时，$\Delta_{i,i}$ 为\Term{diagonal-morphism}{对角态射} $h(X_i)\to h(X_i)\times_{h(X)}h(X_i)$。
\par $\gamma$) 当 $i\ne j$ 时，$\Delta_{i,j}$ 为\Term{morphism}{态射} $\phi_{\widehat C}\to h(X_i)\times_{h(X)}h(X_j)$（$\phi_{\widehat C}$ 表示 $\widehat C$ 的\Term{initial-object}{始对象}）。
\Term{morphism}{态射} $a(\Phi)$ 是\Term{monomorphism}{单态射}当且仅当 $a(\Delta)$ 是\Term{isomorphism}{同构}；由 4.1，$a(\Delta)$ 是\Term{isomorphism}{同构}当且仅当各 $a(\Delta_{i,j})$ 都是\Term{isomorphism}{同构}，即当且仅当有性质 b)、c)。
\end{proof}
\begin{SGAresult}{推论}{4.6.1}
设 $C$ 是一个 $\mathcal U$-\Term{category}{范畴}，$X$ 是 $C$ 的一个\Term{object}{对象}。
\par 1) 设 $T$ 是 $C$ 上的一个 $\mathcal U$-\Term{topology}{拓扑}。$C$ 的\Term{object}{对象} $X$ 经“\Term{associated-sheaf}{伴随层}”\Term{functor}{函子}（对于\Term{topology}{拓扑} $T$）变为 $\widetilde C_T$ 的\Term{initial-object}{始对象}，当且仅当空\Term{sieve}{筛}\Term{covering}{覆盖} $X$。
\par 2) 当 $X$ 是 $C$ 的\Term{strict-initial-object}{严格始对象}时，对于每个比\Term{canonical-topology}{典范拓扑}细的 $\mathcal U$-\Term{topology}{拓扑}，$X$ 的\Term{associated-sheaf}{伴随层}都是\Term{initial-object}{始对象}。
\end{SGAresult}
\begin{proof}
1) 在 4.6 的 2) 中取指标集 $I$ 为空集。
\par 2) 由 1) 和 4.6 的 1)，只须证明空\Term{sieve}{筛}对于\Term{canonical-topology}{典范拓扑}\Term{covering}{覆盖} $X$；即由 2.6，只须证明对 $X$ 上方的每个\Term{object}{对象} $Y\to X$，$Y$ 是 $X$ 的\Term{initial-object}{始对象}。这由定义 4.5 得出。
\end{proof}
\begin{SGAresult}{推论}{4.6.2}
设 $C$ 是一个 $\mathcal U$-\Term{site}{景}，$(s_i:X_i\to X)$，$i\in I$，是 $C$ 中以同一\Term{object}{对象}为靶、\Term{quarrable-morphism}{可作基变换}的一族\Term{morphism}{态射}，满足下列性质：
\par $\alpha$) 族 $s_i$，$i\in I$，是\Term{covering-family}{覆盖族}。
\par $\beta$) 对每个 $i\in I$，$s_i:X_i\to X$ 是\Term{monomorphism}{单态射}。
\SourcePage{149}
\par $\gamma$) 对 $I$ 的每对元素 $(i,j)$，$i\ne j$，$X_i\times_XX_J$ 被空\Term{sieve}{筛}\Term{covering}{覆盖}，即对 $C$ 上的每个\Term{sheaf}{层} $F$，$F(X_i\times_XX_J)$ 是单元素集。
则 $\varepsilon_C(X)$ 是各 $\varepsilon_C(X_i)$，$i\in I$，的和。
\end{SGAresult}
\begin{proof}由 4.6 的 2) 和 4.6.1 的 1) 立即得出。\end{proof}
\begin{SGAresult}{推论}{4.6.3}
设 $C$ 是一个 $\mathcal U$-\Term{category}{范畴}，$(s_i:X_i\to X)$，$i\in I$，是以同一\Term{object}{对象}为靶、\Term{quarrable-morphism}{可作基变换}的一族\Term{morphism}{态射}。假定 $C$ 的\Term{canonical-topology}{典范拓扑}是一个 $\mathcal U$-\Term{topology}{拓扑}。下列条件\Term{equivalence}{等价}：
\par i) 存在 $C$ 上一个比\Term{canonical-topology}{典范拓扑}粗的 $\mathcal U$-\Term{topology}{拓扑} $T$，使 $\varepsilon_T(X)$ 是各 $\varepsilon_T(X_i)$，$i\in I$，的和。
\par ii) \Term{object}{对象} $X$ 是各 $X_i$ 在 $C$ 中的\Term{universal-disjoint-coproduct}{普适不交和}（4.5）。
\end{SGAresult}
\begin{proof}
ii)$\Rightarrow$i)。由于 $X$ 是各 $X_i$，$i\in I$，的\Term{universal-coproduct}{普适和}，族 $s_i$，$i\in I$，对于 $C$ 的\Term{canonical-topology}{典范拓扑}是\Term{covering-family}{覆盖族}（2.6）。故当 $C$ 配备\Term{canonical-topology}{典范拓扑}时，4.6.2 的条件 $\alpha$) 成立。条件 $\beta$) 显然成立，性质 $\gamma$) 则由 4.6 1) 2) 得出，故有 i)。

i)$\Rightarrow$ii)。设 $T$ 是 $C$ 上比\Term{canonical-topology}{典范拓扑}粗的\Term{topology}{拓扑}，且 $\varepsilon_T(X)$ 是各 $\varepsilon_T(X_i)$ 的和。对每个对于 $T$ 的\Term{sheaf}{层} $F$，\Term{canonical}{典范}映射 $\operatorname{Hom}(X,F)\to\prod_i\operatorname{Hom}(X_i,F)$ 是\Term{bijection}{双射}。此外，每个\Term{representable-presheaf}{可表示预层}都是\Term{sheaf}{层}，故立即推出 $X$ 是各 $X_i$，$i\in I$，在 $C$ 中的和。将此用于 $I$ 为空集的情形，可见若 $C$ 的某个\Term{object}{对象}变为 $\widetilde C$ 的\Term{initial-object}{始对象}，则这个\Term{object}{对象}是 $C$ 的\Term{initial-object}{始对象}。\Term{functor}{函子} $\varepsilon_T:C\to\widetilde C$ \Term{fully-faithful}{全忠实}且与\Term{finite-projective-limit}{有限投射极限}交换。于是，4.6.2) 的条件 b) 蕴涵各 $s_i$，$i\in I$，是 $C$ 的\Term{monomorphism}{单态射}，而条件 c) 由上述论证蕴涵各 $X_i\times_XX_j$，$i\ne j$，是 $C$ 的\Term{initial-object}{始对象}。因此，$X$ 是各 $X_i$ 的\Term{disjoint-coproduct}{不交和}。由于 $\varepsilon_T$ 与\Term{fiber-product}{纤维积}交换，这一性质对任意\Term{base-change}{基变换} $Y\to X$ \Term{stable}{稳定}，故 $X$ 是各 $X_i$，$i\in I$，的\Term{universal-disjoint-coproduct}{普适不交和}。
\end{proof}
\begin{SGAresult}{推论}{4.6.4}
设 $C$ 是一个 $\mathcal U$-\Term{category}{范畴}，$X$ 是 $C$ 的\Term{object}{对象}。假定 $C$ 上的\Term{canonical-topology}{典范拓扑}是一个 $\mathcal U$-\Term{topology}{拓扑}。下列性质\Term{equivalence}{等价}：
\par i) 存在 $C$ 上一个比\Term{canonical-topology}{典范拓扑}粗的\Term{topology}{拓扑}，使 $\varepsilon_T(X)$ 是 $\widetilde C_T$ 的\Term{initial-object}{始对象}。
\par ii) $X$ 是 $C$ 的\Term{strict-initial-object}{严格始对象}。
\end{SGAresult}
\begin{proof}在 4.6.3 中取指标集 $I$ 为空集。\end{proof}
\begin{SGAresult}{命题}{4.7}
设 $C$ 是一个 $\mathcal U$-\Term{category}{范畴}，$\widetilde C$ 是 $C$ 对于\Term{canonical-topology}{典范拓扑}的\Term{sheaf}{层}\Term{category}{范畴}，$\varepsilon_C:C\to\widetilde C$ 是\Term{canonical-functor}{典范函子}（4.4.0），$R\rightrightarrows X$ 是 $C$ 中一个
\SourcePage{150}
在 $C$ 中具有\Term{coequalizer}{余等化子} $Y$ 的\Term{equivalence-relation}{等价关系}。设 $\Pi:X\to Y$ 是\Term{canonical}{典范}\Term{morphism}{态射}，并假定它\Term{quarrable-morphism}{可作基变换}。考虑下列性质：
\par i) $\varepsilon_C(Y)$ 是\Term{equivalence-relation}{等价关系} $\varepsilon_C(R)\rightrightarrows\varepsilon_C(X)$ 的商。
\par ii) \Term{equivalence-relation}{等价关系} $R\rightrightarrows X$ \Term{universally-effective}{普适有效}（参见第 7 节）。
有 ii)$\Rightarrow$i)。当 $C$ 有一个比\Term{canonical-topology}{典范拓扑}粗的 $\mathcal U$-\Term{topology}{拓扑}时，有 i)$\Rightarrow$ii)。
\end{SGAresult}
\begin{proof}
ii)$\Rightarrow$i)。\Term{canonical}{典范}\Term{morphism}{态射} $R\to X\times_XX$ 是\Term{isomorphism}{同构}。设 $S\to Y$ 是由 $\pi:X\to Y$ 生成的\Term{sieve}{筛}。对每个\Term{presheaf}{预层} $F$，有\Term{exact}{正合}图（第 I 讲 2.12）
\[\operatorname{Hom}(S,F)\longrightarrow\operatorname{Hom}(X,F)\rightrightarrows\operatorname{Hom}(R,F).\]
因此，只须证明对每个对于\Term{canonical-topology}{典范拓扑}的\Term{sheaf}{层} $F$，映射
\[\operatorname{Hom}(X,F)\longrightarrow\operatorname{Hom}(S,F)\]
是\Term{bijection}{双射}，即只须证明 $S$ 对于\Term{canonical-topology}{典范拓扑}是\Term{covering-sieve}{覆盖筛}，亦即由 $\pi:X\to Y$ 生成的\Term{sieve}{筛}是\Term{covering-sieve}{覆盖筛}。这由定义 2.5 得出。

i)$\Rightarrow$ii)。\Term{functor}{函子} $\varepsilon_C:C\to\widetilde C$ 与\Term{finite-projective-limit}{有限投射极限}交换且\Term{fully-faithful}{全忠实}。而\Term{canonical}{典范}\Term{morphism}{态射} $\varepsilon_C(R)\to J_C(X)\times_{J_C(Y)}J_C(X)$ 是\Term{isomorphism}{同构}（4.3）。所以\Term{canonical}{典范}\Term{morphism}{态射} $R\to X\times_YX$ 是\Term{isomorphism}{同构}。\Term{morphism}{态射} $\varepsilon_C(X)\to\varepsilon_C(Y)$ 是\Term{epimorphism}{满态射}，由此推出（4.4），$\pi:X\to Y$ 是 $C$ 对于\Term{canonical-topology}{典范拓扑}的\Term{covering-morphism}{覆盖态射}。
\end{proof}
\begin{SGAresult}{命题}{4.8}
设 $C$ 是一个 $\mathcal U$-\Term{site}{景}。$C$ 上的\Term{sheaf}{层}\Term{category}{范畴}具有下列性质：
\par a) \Term{finite-projective-limit}{有限投射极限}\Term{representable}{可表示}。
\par b) 以 $\mathcal U$ 的元素为指标的\Term{direct-sum}{直和}\Term{representable}{可表示}，且是不交而\Term{universal}{普适}的（4.5）。
\par c) \Term{equivalence-relation}{等价关系}\Term{universally-effective}{普适有效}。
\end{SGAresult}
这些性质已在 4.1 的 2)、3)，4.3 的 3) 和 4.5.1 中证明。
特意列出这些性质，是因为以后将用它们刻画\Term{equivalence}{等价}于属于 $\mathcal U$ 的\Term{category}{范畴}上的\Term{sheaf}{层}\Term{category}{范畴}的 $\mathcal U$-\Term{category}{范畴}（第 IV 讲 1.2）。

\textbf{注记。}回忆，性质 a) \Term{equivalence}{等价}于：存在\Term{final-object}{终对象}，且\Term{fiber-product}{纤维积}\Term{representable}{可表示}（第 I 讲 2.3.1）。

\textbf{4.9。}设 $C$ 是底\Term{category}{范畴}为 $\mathcal U$-\Term{category}{范畴}的\Term{site}{景}。则 $C$ 上的 $\mathcal U$-\Term{sheaf}{层}\Term{category}{范畴} $\mathscr C$ 满足第 I 讲 7.3 所考虑的条件（(i) 或 (ii)，任选其一），因而由该处，第 I 讲 7.1 中关于 $\mathscr C$ 的\Term{generating-family}{生成族}概念的各种变式相同。于是：
\SourcePage{151}
\begin{SGAresult}{命题}{4.10}
沿用上述记号，考虑\Term{canonical-functor}{典范函子} $\varepsilon:C\to\widetilde C$（4.4.0），设 $G\subset\operatorname{ob}C$ 是 $C$ 中的\Term{topologically-generating-family}{拓扑生成族}（3.0.1）。则 $\widetilde C$ 的\Term{object}{对象}族 $(\varepsilon(X))_{X\in G}$ 是\Term{generating-family}{生成族}。特别地，$C$ 的\Term{object}{对象}族 $(\varepsilon(X))_{X\in\operatorname{ob}C}$ 是\Term{generating-family}{生成族}。
\end{SGAresult}
须证 $\widetilde C$ 中任何满足下述条件的\Term{morphism}{态射} $u:F\to F'$ 都是\Term{isomorphism}{同构}：对每个 $X\in G$，
\[\operatorname{Hom}(\varepsilon(X),F)\longrightarrow\operatorname{Hom}(\varepsilon(X),F')\tag{4.10.1}\]
是\Term{bijection}{双射}。而这个映射等同于映射
\[u(X):F(X)\longrightarrow F'(X),\tag{4.10.2}\]
须证若它对每个 $X\in G$ 是\Term{bijection}{双射}，则它对每个 $X\in\operatorname{ob}C$ 也如此。取由 $G$ 的\Term{object}{对象}组成的 $X$ 的\Term{covering-family}{覆盖族} $(X_i\to X)_{i\in I}$；对每对指标 $(i,j)$，考虑 $\widehat C$ 中所有\Term{morphism}{态射} $X_{ijk}\to X_i\times_XX_j$ 的\Term{set}{集合}，其中 $X_{ijk}\in G$（$k$ 遍历指标集 $I_{i,j}$）。得到如下\Term{exact}{正合}\Term{set}{集合}图之间的\Term{homomorphism}{同态}：
\[\begin{tikzcd}[column sep=large]
F(X)\arrow[r]\arrow[d]&\prod_i F(X_i)\arrow[r,shift left=0.6ex]\arrow[r,shift right=0.6ex]\arrow[d,"\sim"]&\prod_{ijk}F(X_{ijk})\arrow[d,"\sim"]\\
F(X)\arrow[r]&\prod_i F(X_i)\arrow[r,shift left=0.6ex]\arrow[r,shift right=0.6ex]&\prod_{ijk}F(X_{ijk})
\end{tikzcd}\]
按假设，后两个竖\Term{arrow}{箭头}是\Term{bijection}{双射}，故第一个竖\Term{arrow}{箭头}也是\Term{bijection}{双射}。证毕。
\begin{SGAresult}{推论}{4.11}
假定 $C$ 是一个 $\mathcal U$-\Term{site}{景}（3.0.2）。则\Term{category}{范畴} $\widetilde C$ 是 $\mathcal U$-\Term{category}{范畴}，并有一个 $\mathcal U$-小\Term{generating-family}{生成族}。
\end{SGAresult}
事实上，按假设，在 4.10 中可取 $G$ 为小\Term{generating-family}{生成族}，这证明了小\Term{generating-family}{生成族}的存在。另一方面，若 $F,F'\in\operatorname{ob}\mathscr C$，只要知道对每个 $X\in\operatorname{ob}G$ 的\Term{homomorphism}{同态} (4.10.1)，即 (4.10.2)，就知道从 $F$ 到 $F'$ 的\Term{homomorphism}{同态}（第 I 讲 7.1.1）。于是映射
\[\operatorname{Hom}(F,F')\longrightarrow\prod_{X\in G}\operatorname{Hom}(F(X),F'(X))\]
是\Term{injectivity}{单射}。右端是 $\mathcal U$-小的，故左端也是。这证明 $\widetilde C$ 是一个 $\mathcal U$-\Term{category}{范畴}。
\begin{SGAresult}{推论}{4.12}
设 $C$ 是一个 $\mathcal U$-\Term{site}{景}。则对 $\widetilde C$ 的每个\Term{object}{对象} $X$，$X$ 的\Term{subobject}{子对象}\Term{set}{集合}与\Term{quotient-object}{商对象}\Term{set}{集合}都是 $\mathcal U$-小的。
\end{SGAresult}
这分别由第 I 讲 7.4 和第 I 讲 7.5 得出；根据 4.11，可以应用这两条结果。

\SourcePage{152}
\SGAsection{5}{将 $C$ 上的拓扑扩张到 $\widehat C$}
\begin{SGAresult}{命题}{5.1}
设 $C$ 是一个 $\mathcal U$-\Term{site}{景}，$f:H\to K$ 是 $\widehat C$ 的一个\Term{morphism}{态射}。下列条件\Term{equivalence}{等价}：
\par i) 对每个 $X\to K$，$X\in\operatorname{ob}C$，相应\Term{morphism}{态射} $H\times_KX\to X$ 的像是\Term{covering}{覆盖} $X$ 的\Term{sieve}{筛}。
\par ii) \Term{associated-sheaf}{伴随层}之间的\Term{morphism}{态射} $a(f):aH\to aK$ 是 $\widetilde C$ 的\Term{epimorphism}{满态射}。
\par ii bis) 对 $C$ 上的每个\Term{sheaf}{层} $F$，由 $f$ \Term{derivation}{导出}的映射 $\operatorname{Hom}(K,F)\to\operatorname{Hom}(H,F)$ 是\Term{injectivity}{单射}。
\end{SGAresult}
\begin{proof}
ii) 与 ii bis) 显然\Term{equivalence}{等价}。
证明 (i)$\Rightarrow$(ii)。由 4.4，(i) 意味着 $u:a(H\times_KX)\to a(X)$ 是\Term{epimorphism}{满态射}。由于 $a$ 与\Term{finite-projective-limit}{有限投射极限}交换（4.1），有 $a(H\times_KX)\simeq a(H)\times_{a(K)}a(X)$，故所考虑的\Term{morphism}{态射}由 $a(f):a(H)\to a(K)$ 作\Term{base-change}{基变换}得到。由于 $\widetilde C$ 的\Term{epimorphism}{满态射}\Term{universal}{普适}，这证明 (ii)$\Rightarrow$(i)。反过来，由于各 $X_i\to K$ 所成的族是\Term{epimorphic-family}{满态射族}，各 $a(X)\to a(K)$ 在 $\widetilde C$ 中也组成\Term{epimorphic-family}{满态射族}（4.1）。故为证 $a(f):a(H)\to a(K)$ 为\Term{epimorphism}{满态射}，只须在每个上述类型的\Term{base-change}{基变换} $a(X)\to a(K)$ 后验证；这证明 i)$\Rightarrow$ii)。
\end{proof}
\begin{SGAresult}{定义}{5.2}
1) 满足 5.1 中三个\Term{equivalence}{等价}条件的\Term{morphism}{态射} $f:H\to K$，称为\Term{covering-morphism}{覆盖态射}。以同一\Term{object}{对象}为靶的一族\Term{morphism}{态射} $f_i:H_i\to K$，$i\in I$，称为\Term{covering-family}{覆盖族}，如果相应\Term{morphism}{态射} $f:\coprod_iH_i\to K$ 是\Term{covering-morphism}{覆盖态射}。
\par 2) \Term{morphism}{态射} $f:H\to K$ 称为\Term{bicovering-morphism}{双覆盖态射}（bicouvrant），如果它是\Term{covering-morphism}{覆盖态射}，而且\Term{diagonal-morphism}{对角态射} $H\to H\times_KH$ 也是\Term{covering-morphism}{覆盖态射}。以同一\Term{object}{对象}为靶的一族 $f_i:H_i\to K$，$i\in I$，称为\Term{bicovering-family}{双覆盖族}，如果相应\Term{morphism}{态射} $f:\coprod_iH_i\to K$ 是\Term{bicovering-morphism}{双覆盖态射}。
\end{SGAresult}
由 5.1 的条件 i)，族 $(f_i:H_i\to K)$，$i\in I$，是\Term{covering-family}{覆盖族}，意味着对每个\Term{morphism}{态射} $X\to K$，$X\in\operatorname{ob}C$，族 $H_i\times_KX\to X$，$i\in I$，的像是\Term{covering}{覆盖} $X$ 的\Term{sieve}{筛}；或者由 ii)，意味着 $\widetilde C$ 中的\Term{morphism}{态射}族 $a(f_i):aH_i\to aK$ 是\Term{epimorphic-family}{满态射族}（这里用到\Term{functor}{函子} $a$ 与\Term{direct-sum}{直和}交换（4.1））。\footnote{注意，$\widehat C$ 的一个\Term{morphism}{态射}或\Term{morphism}{态射}族为\Term{covering-morphism}{覆盖态射}或\Term{covering-family}{覆盖族}（相应地，\Term{bicovering-morphism}{双覆盖态射}或\Term{bicovering-family}{双覆盖族}）的性质对\Term{base-change}{基变换}\Term{stable}{稳定}。}
\begin{SGAresult}{命题}{5.3}
设 $C$ 是一个 $\mathcal U$-\Term{site}{景}，$f:H\to K$ 是 $\widehat C$ 的\Term{morphism}{态射}。下列条件\Term{equivalence}{等价}：
\par i) \Term{morphism}{态射} $f$ 是\Term{bicovering-morphism}{双覆盖态射}（5.2）。
\SourcePage{153}
\par i bis) \Term{morphism}{态射} $f$ 是\Term{covering-morphism}{覆盖态射}，且对 $C$ 的每个\Term{object}{对象} $X$ 及满足 $fu=fv$ 的每对\Term{morphism}{态射} $u,v:X\rightrightarrows H$，$(u,v)$ 的\Term{equalizer}{等化子}是\Term{covering}{覆盖} $X$ 的\Term{sieve}{筛}。
\par ii) \Term{morphism}{态射} $a(f):aH\to aK$ 是 $\widetilde C$ 的\Term{isomorphism}{同构}。
\par ii bis) 对 $C$ 上的每个\Term{sheaf}{层} $F$，映射 $\operatorname{Hom}(K,F)\to\operatorname{Hom}(H,F)$ 是\Term{bijection}{双射}。
\end{SGAresult}
\begin{proof}
\Term{equivalence}{等价} i)$\Leftrightarrow$i bis) 由对\Term{diagonal-morphism}{对角态射} $H\to H\times_KH$ 应用 5.1 的条件 i) 得出；\Term{equivalence}{等价} ii)$\Leftrightarrow$ii bis) 是\Term{trivial-bundle}{平凡}的。
\par i)$\Rightarrow$ii)。\Term{morphism}{态射} $a(f):aH\to aK$ 是\Term{epimorphism}{满态射}（5.1）。由于\Term{functor}{函子} $a$ 与\Term{fiber-product}{纤维积}的构造交换（4.1），\Term{diagonal-morphism}{对角态射} $aH\to aH\times_{aK}aH$ 是\Term{epimorphism}{满态射}（5.1）。\Term{diagonal-morphism}{对角态射}总是\Term{monomorphism}{单态射}，故为\Term{isomorphism}{同构}（4.2）。所以\Term{morphism}{态射} $a(f)$ 是\Term{monomorphism}{单态射}，从而是\Term{isomorphism}{同构}（4.2）。
\par ii)$\Rightarrow$i)。由于\Term{functor}{函子} $a$ 与\Term{fiber-product}{纤维积}的构造交换，\Term{morphism}{态射} $f$ 和\Term{diagonal-morphism}{对角态射} $H\to H\times_KH$ 经 $a$ 变为\Term{isomorphism}{同构}。特别地，它们经 $a$ 变为\Term{epimorphism}{满态射}，故都是\Term{covering-morphism}{覆盖态射}。
\end{proof}
\textbf{5.3.1。}由 5.3 和 $a$ 与\Term{direct-sum}{直和}交换可知，$\widehat C$ 中的一族\Term{morphism}{态射} $(f_i:H_i\to K)$，$i\in I$，是\Term{bicovering-family}{双覆盖族}，当且仅当各 $f_i$ 在\Term{associated-sheaf}{伴随层}上诱导从\Term{direct-sum}{直和} $\coprod_i aH_i$ 到 $aK$ 的\Term{isomorphism}{同构}；亦即，当且仅当对每个\Term{sheaf}{层} $F$，映射
\[\operatorname{Hom}(K,F)\longrightarrow\prod_i\operatorname{Hom}(H_i,F)\]
是\Term{bijection}{双射}。
\begin{SGAresult}{命题}{5.4}
设 $C$ 是一个 $\mathcal U$-\Term{site}{景}。$\widehat C$ 上存在一个\Term{topology}{拓扑}（显然唯一），使 $\widehat C$ 中以同一\Term{object}{对象}为靶的\Term{arrow}{箭头}族 $H_i\to K$ 对于这个\Term{topology}{拓扑}是\Term{covering-family}{覆盖族}，当且仅当它在 5.2 的意义下是\Term{covering-family}{覆盖族}。它也是 $\widehat C$ 上具有下列性质的\Term{topology}{拓扑} $T$ 中最粗的一个：
\par a) $T$ 比 $\widehat C$ 的\Term{canonical-topology}{典范拓扑}细（即 $\widehat C$ 的每个\Term{epimorphic-family}{满态射族}对于 $T$ 都是\Term{covering-family}{覆盖族}）。
\par b) $C$ 中的每个\Term{covering-family}{覆盖族}在 $\widehat C$ 中也是\Term{covering-family}{覆盖族}。
\end{SGAresult}
\begin{proof}
我们只给出提示。读者可自行证明，5.2 意义下的\Term{covering-family}{覆盖族}是 $\widehat C$ 上某个\Term{topology}{拓扑} $T_C$ 的\Term{covering-family}{覆盖族}（用到 4.1）。$\widehat C$ 的\Term{canonical-topology}{典范拓扑}的\Term{covering-family}{覆盖族}就是 $\widehat C$ 的\Term{epimorphic-family}{满态射族}（2.6 和第 I 讲 3.1）。由于 $a$ 与\Term{inductive-limit}{归纳极限}交换（4.1），\Term{topology}{拓扑} $T_C$ 比 $\widehat C$ 的\Term{canonical-topology}{典范拓扑}细。此外，$C$ 的\Term{covering-family}{覆盖族}是 $T_C$ 的\Term{covering-family}{覆盖族}（5.1）。故若 $T'$ 表示 $\widehat C$ 上满足性质 a)、
\SourcePage{154}
b) 的最\Term{coarse-topology}{粗拓扑}，则 $T_C$ 比 $T'$ 细。设 $(f_i:H_i\to K)$，$i\in I$，是 $T_C$ 的\Term{covering-family}{覆盖族}。证明 $(f_i,i\in I)$ 是 $T'$ 的\Term{covering-family}{覆盖族}。设 $s_i:H_i\to H=\coprod_iH_i$ 为\Term{canonical}{典范}\Term{monomorphism}{单态射}，$f=(f_i,i\in I):H=\coprod_iH_i\to K$ 为各 $f_i$ 定义的\Term{morphism}{态射}。各 $s_i$ 所成的族对于 $T'$ 是\Term{covering-family}{覆盖族}。因此，由\Term{topology}{拓扑}公理 (T 2)，为证 $(f_i,i\in I)$ 是 $T'$ 的\Term{covering-family}{覆盖族}，只须证明\Term{morphism}{态射} $f:H\to K$ 对于 $T'$ 是\Term{covering-morphism}{覆盖态射}。存在 $\widehat C$ 的\Term{epimorphic-family}{满态射族} $u_\lambda:X_\lambda\to K$，$\lambda\in\Lambda$，$X_\lambda\in\operatorname{ob}C$（第 I 讲 3.4）。故由\Term{topology}{拓扑}公理 (T 2)，只须证明对每个 $\lambda\in\Lambda$，\Term{morphism}{态射} $\mathrm{pr}_2:H\times_KX_\lambda\to X_\lambda$ 对于 $T'$ 是\Term{covering-morphism}{覆盖态射}。再取 $\widehat C$ 的\Term{epimorphic-family}{满态射族} $v_j:Y_j\to H\times_KX_\lambda$，$j\in J$，$Y_j\in\operatorname{ob}C$。族 $(\mathrm{pr}_2\circ v_j,j\in J)$ 对于 $T_C$ 是\Term{covering-family}{覆盖族}，故是 $C$ 的\Term{covering-family}{覆盖族}（5.1），从而是 $T'$ 的\Term{covering-family}{覆盖族}。所以，由 $\mathrm{pr}_2:H\times_KX\to X$ 生成的\Term{sieve}{筛}包含 $T'$ 的一个\Term{covering-sieve}{覆盖筛}，因而对于 $T'$ 是\Term{covering-sieve}{覆盖筛}，所以 $\mathrm{pr}_2:H\times_KX\to X$ 是\Term{covering-morphism}{覆盖态射}。
\end{proof}
\textbf{注记 5.4.1。}5.4 的证明实际上表明，$\widehat C$ 上每个比 $\widehat C$ 的\Term{canonical-topology}{典范拓扑}细的\Term{topology}{拓扑} $T'$，都是 $\widehat C$ 上满足下列性质的\Term{topology}{拓扑} $T$ 中最粗的：
\par a) $T$ 比 $\widehat C$ 的\Term{canonical-topology}{典范拓扑}细。
\par b) 对于 $T'$ 的每个形如 $u_i:X_i\to X$，$i\in I$，其中 $X$ 和各 $X_i$ 都是 $C$ 的\Term{object}{对象}的\Term{covering-family}{覆盖族}，对于 $T$ 也是\Term{covering-family}{覆盖族}。
特别地，$\widehat C$ 上每个比\Term{canonical-topology}{典范拓扑}细的\Term{topology}{拓扑} $T'$，完全由对于 $T'$ 是\Term{covering-family}{覆盖族}的那些族 $u_i:X_i\to X$，$i\in I$，$X_i,X\in\operatorname{ob}C$ 唯一确定。

\textbf{注记 5.4.2。}容易证明，对 $\widehat C$ 上每个比\Term{canonical-topology}{典范拓扑}细的\Term{topology}{拓扑} $T'$，$C$ 中以同一\Term{object}{对象}为靶、且对于 $T'$ 是\Term{covering-family}{覆盖族}的所有\Term{morphism}{态射}族 $(X_i\to X)$，$i\in I$，恰为 $C$ 上某个\Term{topology}{拓扑}的\Term{covering-family}{覆盖族}。因此，5.4 和 5.4.1 在 $C$ 上的\Term{topology}{拓扑}与 $\widehat C$ 上比\Term{canonical-topology}{典范拓扑}细的\Term{topology}{拓扑}之间建立一一\Term{correspondence}{对应}。

\textbf{5.5.0。}设 $C$ 是一个小\Term{category}{范畴}。以 $\operatorname{Caf}$ 表示 $\widehat C$ 的严格\Term{full-subcategory}{满子范畴}（即每个与\Term{subcategory}{子范畴}的\Term{object}{对象}\Term{isomorphism}{同构}的\Term{object}{对象}也属于该\Term{subcategory}{子范畴}）中，其\Term{inclusion-functor}{包含函子}有一个与\Term{finite-projective-limit}{有限投射极限}交换的\Term{left-adjoint}{左伴随}的那些\Term{subcategory}{子范畴}的\Term{set}{集合}。再以 $\mathscr T$ 表示 $C$ 上的\Term{topology}{拓扑}\Term{set}{集合}。定理 3.4 定义一个映射
\[\Phi:\mathscr T\longrightarrow\operatorname{Caf}.\]
下面定义反方向的映射。为此，须对每个元素
\[e=\left(\begin{tikzcd}C'\arrow[r,shift right=0.6ex,"i'"']&\widehat C\arrow[l,shift right=0.6ex,"a'"']\end{tikzcd},\quad a'\text{ 为 }i'\text{ 的左伴随}\right)\]
配上 $C$ 的一个\Term{topology}{拓扑} $T_e$。
\SourcePage{155}
对 $C$ 的每个\Term{object}{对象} $X$，定义 $J_e(X)$ 为那些\Term{embedding}{嵌入}\Term{morphism}{态射}经 $a'$ 变为\Term{isomorphism}{同构}的 $X$ 的\Term{subobject}{子对象}组成的\Term{set}{集合}。利用对 $a'$ 的假设，立即可知这定义了 $C$ 上的一个\Term{topology}{拓扑} $T_e$。于是定义了一个映射
\[\Psi:\operatorname{Caf}\longrightarrow\mathscr T.\]
有如下结果（归功于 J. Giraud）：
\begin{SGAresult}{定理}{5.5}
映射 $\Phi$ 是\Term{bijection}{双射}，$\Psi$ 是它的逆映射。
\end{SGAresult}
\begin{proof}
映射 $\Psi\circ\Phi$ 是恒等映射，这由 5.1 立即得出。映射 $\Phi\circ\Psi$ 也是恒等映射。事实上，设
\[\left(\begin{tikzcd}C'\arrow[r,shift right=0.6ex,"i'"']&\widehat C\arrow[l,shift right=0.6ex,"a'"']\end{tikzcd}\right)\]
是 $\operatorname{Caf}$ 的元素，$T_e$ 是由 $\Psi$ \Term{correspondence}{对应}于它的\Term{topology}{拓扑}，$C_e$ 是对于 $T_e$ 的\Term{sheaf}{层}\Term{category}{范畴}。利用\Term{topology}{拓扑} $T_e$ 的定义与\Term{bicovering-morphism}{双覆盖态射}的定义（5.2），可证明下列断言\Term{equivalence}{等价}：
\par i) $\widehat C$ 的\Term{morphism}{态射} $u$ 对于\Term{topology}{拓扑} $T_e$ 是\Term{bicovering-morphism}{双覆盖态射}。
\par ii) $\widehat C$ 的\Term{morphism}{态射} $u$ 经 $a'$ 变为\Term{isomorphism}{同构}。
显然，$C'$ 是 $C_e$ 的\Term{full-subcategory}{满子范畴}，故只须证明对于\Term{topology}{拓扑} $T_e$ 的每个\Term{sheaf}{层} $F$ 都是 $C'$ 的\Term{object}{对象}。由上述\Term{equivalence}{等价}，\Term{morphism}{态射} $F\to i'\circ a'(F)$ 是\Term{bicovering-morphism}{双覆盖态射}，而它的源与靶都是\Term{sheaf}{层}，所以由 5.3 的 ii)，$F\to i'\circ a'(F)$ 是\Term{isomorphism}{同构}。
\end{proof}

\SGAsection{6}{取值于范畴的层}
\textbf{6.0。}设 $C$、$D$ 是两个\Term{category}{范畴}。从 $C$ 到 $D$ 的\Term{contravariant-functor}{反变函子} $F:C^\circ\to D$，称为 $C$ 上取值于 $D$ 的\Term{presheaf}{预层}。
\begin{SGAresult}{定义}{6.1}
设 $C$ 是一个\Term{site}{景}，$D$ 是一个\Term{category}{范畴}。\Term{presheaf}{预层} $F:C^\circ\to D$ 称为 $C$ 上取值于 $D$ 的\Term{sheaf}{层}（简称取值于 $D$ 的\Term{sheaf}{层}），如果对 $D$ 的每个\Term{object}{对象} $S$，\Term{presheaf-of-sets}{集合预层}
\[X\longmapsto\operatorname{Hom}_D(S,F(X)),\qquad X\in\operatorname{ob}(C)\]
是\Term{sheaf}{层}。$\mathscr Hom(C^\circ,D)$ 中由 $C$ 上取值于 $D$ 的\Term{sheaf}{层}构成的\Term{full-subcategory}{满子范畴}，记为 $\mathscr Hom^{\sim}(C^\circ,D)$。
\end{SGAresult}
\textbf{注记 6.2。}1) 条件 6.1 意味着对 $C$ 的每个\Term{object}{对象} $X$、\Term{covering}{覆盖} $X$ 的每个\Term{sieve}{筛} $R$ 以及 $D$ 的每个\Term{object}{对象} $S$，有\Term{isomorphism}{同构}（第 I 讲 3.5 和 2.1）：
\[\operatorname{Hom}_D(S,F(X))\xrightarrow{\sim}\varprojlim_{C/R}F(\,\cdot\,).\]
当 $D$ 是 $\mathcal U$-\Term{set}{集合}\Term{category}{范畴}时，这就给出\Term{sheaf-of-sets}{集合层}的定义（同上）。
\SourcePage{156}
2) 设 $D'$ 是一个 $\mathcal U$-\Term{category}{范畴}，$G:D\to D'$ 是与\Term{projective-limit}{投射极限}交换的\Term{functor}{函子}。由复合，\Term{functor}{函子} $G$ 把取值于 $D$ 的\Term{sheaf}{层}变为取值于 $D'$ 的\Term{sheaf}{层}。

\textbf{6.3.0。}设 $\gamma$ 是由\Term{finite-projective-limit}{有限投射极限}定义的一类\Term{algebraic-structure}{代数结构}（第 I 讲 2.9），$\mathcal U\text{-}\gamma$ 是属于 $\mathcal U$ 的 $\gamma$-\Term{set}{集合}的\Term{category}{范畴}。以 $\mathrm{esj}:\mathcal U\text{-}\gamma\to\mathcal U\text{-}\mathrm{Ens}$ 表示“\Term{underlying-set}{底集}合”\Term{functor}{函子}。\Term{functor}{函子} $\mathrm{esj}$ 与\Term{projective-limit}{投射极限}交换，因此由上一条注记，通过复合定义\Term{functor}{函子}
\[\begin{split}
\widehat{\mathrm{esj}}:\mathscr Hom(C^\circ,\mathcal U\text{-}\gamma)&\longrightarrow\widehat C,\\
\widetilde{\mathrm{esj}}:\mathscr Hom^{\sim}(C^\circ,\mathcal U\text{-}\gamma)&\longrightarrow\widetilde C.
\end{split}\]
\Term{functor}{函子} $\widetilde{\mathrm{esj}}$ 称为“底\Term{sheaf-of-sets}{集合层}”\Term{functor}{函子}。事实上，它\Term{canonical}{典范}地经过 $\widetilde C$ 的 $\gamma$-\Term{object}{对象}\Term{category}{范畴} $\gamma\text{-}\widetilde C$ \Term{resolution}{分解}；仍以 $\widetilde{\mathrm{esj}}$ 表示所得到的\Term{functor}{函子}
\[\mathscr Hom^{\sim}(C^\circ,\mathcal U\text{-}\gamma)\longrightarrow\gamma\text{-}\widetilde C,\]
则有：
\begin{SGAresult}{命题}{6.3.1}
\Term{functor}{函子} $\widetilde{\mathrm{esj}}$ 在下列三个\Term{category}{范畴}之间建立\Term{equivalence-of-categories}{范畴等价}：$C$ 上取值于 $\mathcal U\text{-}\gamma$ 的\Term{sheaf}{层}\Term{category}{范畴}、\Term{sheaf-of-sets}{集合层}\Term{category}{范畴}的 $\gamma$-\Term{object}{对象}\Term{category}{范畴}，以及\Term{presheaf-of-sets}{集合预层}\Term{category}{范畴}中底\Term{presheaf}{预层}为\Term{sheaf}{层}的 $\gamma$-\Term{object}{对象}\Term{category}{范畴}。
\end{SGAresult}
\begin{proof}证明主要利用第 I 讲 3.2，以及在\Term{presheaf}{预层}\Term{category}{范畴}中\Term{sheaf}{层}的\Term{finite-projective-limit}{有限投射极限}仍是\Term{sheaf}{层}这一事实（4.1）。\end{proof}
\textbf{6.3.2。}上述命题说明，可以不甚严格地把取值于 $\mathcal U\text{-}\gamma$ 的\Term{sheaf}{层}与 $\widetilde C$ 中相应的 $\gamma$-\Term{object}{对象}等同起来。今后将一贯采用这一说法。
我们使用经典术语：\Term{group-sheaf}{群层}、\Term{sheaf-of-commutative-groups}{交换群层}（通常称为\Term{abelian-sheaf}{阿贝尔层}）、\Term{ring-sheaf}{环层}、$A$-\Term{module-sheaf}{模层}……

\textbf{6.3.3。}以
\[\widetilde C_\gamma\quad(\text{相应地，}\widehat C_\gamma)\]
表示 $\widetilde C$（相应地，$\widehat C$）的 $\gamma$-\Term{object}{对象}\Term{category}{范畴}。当 $\gamma$ 为“$A$-模”这一类结构时，也记为 $\widetilde C_A$；当 $A=\mathbf Z$ 时，简称为 $\widetilde C_{\mathrm{ab}}$。
假设 $C$ 是一个 $\mathcal U$-\Term{site}{景}。“\Term{associated-sheaf}{伴随层}”\Term{functor}{函子}\Term{left-exact}{左正合}（4.1），因而：
\SourcePage{157}
\begin{SGAresult}{命题}{6.4}
\Term{inclusion-functor}{包含函子} $i_\gamma:\widetilde C_\gamma\to\widehat C_\gamma$ 有一个\Term{left-exact}{左正合}的\Term{left-adjoint}{左伴随} $a_\gamma$，即有\Term{isomorphism}{同构}
\[\operatorname{Hom}_{\widetilde C_\gamma}(a_\gamma\,\cdot\,,\,\cdot\,)
\xrightarrow{\sim}\operatorname{Hom}_{\widehat C_\gamma}(\,\cdot\,,i_\gamma\,\cdot\,).\]
设 $X$ 是 $\widehat C_\gamma$ 的一个\Term{object}{对象}。$a_\gamma(X)$ 的底\Term{sheaf-of-sets}{集合层}\Term{canonical}{典范}\Term{isomorphism}{同构}于 $X$ 的底\Term{presheaf-of-sets}{集合预层}的\Term{adjunction}{伴随}\Term{sheaf-of-sets}{集合层}。\Term{adjunction-morphism}{伴随态射} $\mathrm{id}\to i_\gamma a_\gamma$ 的底\Term{presheaf-of-sets}{集合预层}\Term{morphism}{态射}，等同于对底\Term{presheaf-of-sets}{集合预层}应用\Term{adjunction-morphism}{伴随态射}所得的\Term{morphism}{态射}。
\end{SGAresult}
\begin{SGAresult}{命题}{6.5}
假定\Term{functor}{函子} $\mathrm{esj}:\gamma\text{-}\mathcal U\to\mathcal U\text{-}\mathrm{Ens}$ 有一个\Term{left-adjoint}{左伴随}\footnote{可以证明，这样的\Term{adjunction}{伴随}总是存在 [C.F.]：\Term{free-group}{自由群}、自由\Term{abelian-group}{阿贝尔群}、自由 $A$-模等。}：
\[\operatorname{Hom}_{\mathcal U\text{-}\mathrm{Ens}}(\,\cdot\,,\mathrm{esj}(\,\cdot\,))
\xrightarrow{\sim}\operatorname{Hom}_{\gamma\text{-}\mathcal U}(\mathrm{Lib}(\,\cdot\,),\,\cdot\,).\]
\Term{functor}{函子} $\widetilde{\mathrm{esj}}:\widetilde C_\gamma\to\widetilde C$（相应地，$\widehat{\mathrm{esj}}:\widehat C_\gamma\to\widehat C$）有一个\Term{left-adjoint}{左伴随} $\widetilde{\mathrm{Lib}}$（相应地，$\widehat{\mathrm{Lib}}$），并有\Term{canonical}{典范}\Term{isomorphism}{同构}
\[\widetilde{\mathrm{Lib}}=a_\gamma\widehat{\mathrm{Lib}}\,i.\]
\end{SGAresult}
\begin{proof}证明是形式的，留给读者。\end{proof}
\begin{SGAresult}{推论}{6.6}
设 $(X_i)$，$i\in I$，是 $\widetilde C$ 的一个\Term{generating-family}{生成族}（4.9）。则族 $\widetilde{\mathrm{Lib}}(X_i)$ 是 $\widetilde C_\gamma$ 的\Term{generating-family}{生成族}。
\end{SGAresult}
\begin{proof}只须注意，\Term{functor}{函子} $\widetilde{\mathrm{esj}}$ 与\Term{projective-limit}{投射极限}交换且\Term{conservative}{保守}（$\widetilde{\mathrm{esj}}(u)$ 是\Term{isomorphism}{同构} $\Leftrightarrow u$ 是\Term{isomorphism}{同构}），其余证明是形式的。\end{proof}
\begin{SGAresult}{命题}{6.7}
设 $A$ 是一个 $\mathcal U$-\Term{ring-sheaf}{环层}，或者一个小环，$\widetilde C_A$（相应地，$\widehat C_A$）是 $\mathcal U$-\Term{site}{景} $C$ 上的幺 $A$-\Term{module-sheaf}{模层}（相应地，\Term{presheaf}{预层}）\Term{category}{范畴}（6.3.3）。则 $\widetilde C_A$ 是一个满足 [Tohoku] 的公理 (AB 5)（“存在\Term{left-exact}{左正合}的\Term{filtered-inductive-limit}{滤过归纳极限}”）和 (AB 3)$^*$（“存在无限积”）的阿贝尔 $\mathcal U$-\Term{category}{范畴}。它有一个以 $\mathcal U$ 的元素为指标的\Term{family-of-generators}{生成元族}。
\end{SGAresult}
\begin{proof}
显然，\Term{category}{范畴} $\widehat C_A$ 是满足公理 (AB 3)、(AB 4)、(AB 5) 的\Term{abelian-category}{阿贝尔范畴}（第 I 讲 2.8 和第 I 讲 3.3）。由 6.4 推出，\Term{category}{范畴} $\widetilde C_A$ 是核与\Term{cokernel}{余核}都\Term{representable}{可表示}的\Term{additive-category}{可加范畴}。更确切地，设 $i_A$、$a_A$ 分别是到\Term{presheaf}{预层}的\Term{inclusion-functor}{包含函子}和\Term{associated-sheaf-functor}{伴随层函子}，$u:X\to Y$ 是 $\widetilde C_A$ 的\Term{morphism}{态射}。\Term{functor}{函子} $a_A$ \Term{left-exact}{左正合}且与\Term{inductive-limit}{归纳极限}交换，\Term{functor}{函子} $i_A$ 与
\SourcePage{158}
\Term{projective-limit}{投射极限}交换。因此\Term{functor}{函子} $i_Aa_A:\widehat C\to\widehat C_A$ 是\Term{left-exact}{左正合}\Term{additive-functor}{可加函子}。由此得到\Term{canonical}{典范}\Term{isomorphism}{同构}
\[i_a(\operatorname{Ker}(u))\xrightarrow{\sim}\operatorname{Ker}(i_A(u)),\tag{*}\]
\[\operatorname{coker}(u)\xrightarrow{\sim}a_A\operatorname{coker}(i_a(u)).\tag{**}\]
所以有\Term{canonical}{典范}\Term{isomorphism}{同构}
\[\operatorname{coim}(u)\xrightarrow{\sim}a_A\operatorname{coim}(i_a(u)).\tag{***}\]
证明\Term{canonical}{典范}\Term{morphism}{态射} $\operatorname{coim}(u)\to\operatorname{im}(u)$ 为\Term{isomorphism}{同构}。由 (*)，只须证明序列
\[0\longrightarrow i_A(\operatorname{coim}(u))\longrightarrow i_A(Y)\longrightarrow i_A(\operatorname{coker}(u))\]
\Term{exact}{正合}。利用\Term{isomorphism}{同构} (**)、(***)，并注意 $a_Ai_A(Y)$ \Term{isomorphism}{同构}于 $Y$，可见这个序列\Term{isomorphism}{同构}于序列
\[0\longrightarrow i_Aa_A(\operatorname{coim}(i_A(u)))\longrightarrow i_Aa_Ai_A(Y)
\longrightarrow i_Aa_A(\operatorname{coker}(i_A(u))),\]
后者正是序列
\[0\longrightarrow\operatorname{coim}(i_A(u))\longrightarrow i_A(Y)\longrightarrow\operatorname{coker}(i_A(u))\]
经\Term{functor}{函子} $i_Aa_A$ 的像。这个序列\Term{exact}{正合}，而\Term{functor}{函子} $i_Aa_A$ \Term{left-exact}{左正合}，所以\Term{category}{范畴} $\widetilde C_A$ 是\Term{abelian-category}{阿贝尔范畴}。\Term{functor}{函子} $a_A:\widehat C_A\to\widetilde C_A$ \Term{left-exact}{左正合}且与任意\Term{inductive-limit}{归纳极限}交换，因此它可加、\Term{exact}{正合}，且与\Term{inductive-limit}{归纳极限}交换。\Term{category}{范畴} $\widehat C_A$ 满足公理 (AB 5)，故\Term{category}{范畴} $\widetilde C_A$ 也满足。最后，显然在\Term{category}{范畴} $\widetilde C_A$ 中，以 $\mathcal U$ 的元素为指标的积\Term{representable}{可表示}，故公理 (AB 3)$^*$ 成立。这就证明了第一个断言。

对每个属于 $\mathcal U$ 的环 $A$，\Term{functor}{函子}
\[\mathrm{esj}:\mathcal U\text{-}A\text{-}\mathrm{module}\longrightarrow\mathcal U\text{-}\mathrm{Ens}\]
有一个\Term{left-adjoint}{左伴随} $X_A$（所生成的自由 $A$-模）。故应用 4.10 和 6.6 即完成证明。
\end{proof}
\textbf{记号 6.8。}设 $H$ 是一个\Term{sheaf-of-sets}{集合层}。\Term{adjunction}{伴随}于\Term{presheaf}{预层}（参见记号 6.5）
\[X\longmapsto\mathrm{Lib}_A H(X),\qquad X\in\operatorname{ob}C\]
的 $A$-\Term{module-sheaf}{模层}记为 $A_H$。当 $H=a(X)$（由 $X$ 表示的\Term{presheaf}{预层}的\Term{associated-sheaf}{伴随层}）时，有时也不甚严格地简称为 $A_X$。6.7 的证明表明，$A$-\Term{module-sheaf}{模层}族 $\mathscr A_X$，$X\in\operatorname{ob}(C)$，是 $A$-\Term{module-sheaf}{模层}\Term{category}{范畴}的一个以 $\mathcal U$ 的元素为指标的\Term{generating-family}{生成族}。
\SourcePage{159}
\textbf{注记 6.9。}根据 [Tôhoku]，6.7 表明，当 $A$ 是 $\mathcal U$ 的一个元素时，\Term{category}{范畴} $\widetilde C_A$ 有足够多的\Term{injective-object}{内射对象}。另一方面，已知在 $\widetilde C_A$ 中无限积不一定\Term{exact}{正合}，所以\Term{category}{范畴} $\widetilde C_A$ 一般没有足够多的\Term{projective-object}{投射对象} [Roos]。

\bigskip\noindent\textbf{参考文献}
\begin{enumerate}
\item[{[1]}] M. Artin；《Grothendieck \Term{topology}{拓扑}》（Grothendieck's topologies）。
\item[{[2]}] M. Demazure：《\Term{algebra}{代数}几何讨论班 III，第 IV 讲》（Séminaire de Géométrie Algébrique III, Exposé IV），Lecteur Notes，Springer-Verlag。
\item[{[3]}] A. Grothendieck：《\Term{homology}{同调}\Term{algebra}{代数}的若干问题》（Sur quelques points d'Algèbre Homologique），Tohoku Math. Journal。
\item[{[4]}] J. E. Roos：CR。
\end{enumerate}
\SourcePage{160}
% 原 PDF 第 160 页为空白页。
